\documentclass[a4paper,11pt]{book}
\usepackage{listings}
\usepackage[utf8]{inputenc}
\usepackage{titlesec}
\usepackage{fancyhdr}
\usepackage[spanish,es-tabla]{babel}
\usepackage[hidelinks]{hyperref}
\usepackage{xcolor}
\usepackage{pdfpages}
\usepackage{amssymb}

% Información reutilizable
\newcommand{\asunto}{Trabajo de Fin de Máster}
\newcommand{\titulo}{Datos Abiertos en la Universidad}
\newcommand{\tituloEng}{Open Data at University}
\newcommand{\master}{Máster Universitario en Ingeniería Informática}
\newcommand{\keywordsEs}{datos abiertos, web semántica, datos enlazados, CKAN, RDF, OWL, SPARQL}
\newcommand{\keywordsEn}{open data, semantic web, CKAN, RDF, OWL, SPARQL}
\newcommand{\autor}{Germán Martínez Maldonado}
\newcommand{\email}{germaaan@gmail.com}
\newcommand{\tutor}{Juan Julián Merelo Guervós}
\newcommand{\escuela}{Escuela Técnica Superior de Ingenierías Informática y de Telecomunicación}
\newcommand{\universidad}{Universidad de Granada}
\newcommand{\ciudad}{Granada}
\newcommand{\vers}{Versión 0.1}

% Información archivo
\hypersetup{
	pdfauthor = {\autor\ (\email)},
	pdftitle = {\titulo},
	pdfsubject = {\asunto},
	pdfkeywords = {\keywordsEs},
	pdfcreator = {LaTeX con el paquete texlive},
	pdfproducer = {pdflatex}
}

% Estilo de cabeceras
\pagestyle{fancy}
\fancyhf{}
\fancyhead[LO]{\leftmark}
\fancyhead[RE]{\rightmark}
\fancyhead[RO,LE]{\textbf{\thepage}}
\setlength{\headheight}{1.5\headheight}

% Redefinición de comandos
\renewcommand{\lstlistingname}{Fragmento de código}
\renewcommand{\lstlistlistingname}{Índice de fragmentos de código}
\renewcommand{\chaptermark}[1]{\markboth{\textbf{#1}}{}}
\renewcommand{\sectionmark}[1]{\markright{\textbf{\thesection. #1}}}
\newcommand\tab[1][0.5cm]{\hspace*{#1}}

% Definición de colores
\definecolor{gray97}{gray}{.97}
\definecolor{gray75}{gray}{.75}
\definecolor{gray45}{gray}{.45}
\definecolor{gray30}{gray}{.94}
\definecolor{lightgray}{rgb}{.9,.9,.9}
\definecolor{darkgray}{rgb}{.4,.4,.4}
\definecolor{purple}{rgb}{0.65, 0.12, 0.82}
\definecolor{background}{HTML}{EEEEEE}
\definecolor{delim}{RGB}{20,105,176}
\colorlet{punct}{red!60!black}
\colorlet{numb}{magenta!60!black}

% Listados
\lstset{
	aboveskip=0.5cm,
	backgroundcolor=\color{gray97},
	basicstyle=\scriptsize\ttfamily,
	breaklines=true,
	commentstyle=\color{gray45},
	frame=Ltb,
	framerule=0.5pt,
	framesep=0pt,
	framexbottommargin=3pt,
	framexleftmargin=0.1cm,
	framextopmargin=3pt,
	keywordstyle=\bfseries,
	numberfirstline = false,
	numbers=left,
	numbersep=6pt,
	numberstyle=\tiny,
	rulesep=.4pt,
	rulesepcolor=\color{black},
	showstringspaces = false,
	stringstyle=\ttfamily,
	literate={á}{{\'a}}1
	         {é}{{\'e}}1
	         {í}{{\'i}}1
	         {ó}{{\'o}}1
	         {ú}{{\'u}}1
	         {ñ}{{\~n}}1
}
 
% Minimizar fragmentado de listados
\lstnewenvironment{listing}[1][]
	{\lstset{#1}\pagebreak[0]}{\pagebreak[0]}

% Para que las páginas en blanco no tengan cabecera
\makeatletter
\def\clearpage{%
  \ifvmode
    \ifnum \@dbltopnum =\m@ne
      \ifdim \pagetotal <\topskip
        \hbox{}
      \fi
    \fi
  \fi
  \newpage
  \thispagestyle{empty}
  \write\m@ne{}
  \vbox{}
  \penalty -\@Mi
}
\makeatother

\begin{document}
\begin{titlepage}

\newlength{\centeroffset}
\setlength{\centeroffset}{-0.5\oddsidemargin}
\addtolength{\centeroffset}{0.5\evensidemargin}

\noindent\hspace*{\centeroffset}\begin{minipage}{\textwidth}

\centering
\includegraphics[width=0.9\textwidth]{../images/logo_ugr.png}\\[1.4cm]

\textsc{\Large\asunto\\[0.2cm]}
\textsc{\master}\\[1cm]

{\Huge\bfseries \titulo\\}
\noindent\rule[-1ex]{\textwidth}{3pt}\\[3.5ex]
\end{minipage}

\vspace{2cm}
\noindent\hspace*{\centeroffset}\begin{minipage}{\textwidth}
\centering

\textbf{Autor}\\ {\autor}\\[2.5ex]
\textbf{Tutor}\\ {\tutor}\\[2cm]
\includegraphics[width=0.3\textwidth]{../images/logo_etsiit.png}\\[0.1cm]
\textsc{\escuela}\\
\textsc{---}\\
\ciudad, \today\\
\includegraphics[width=0.3\textwidth]{../images/CC-SA-logo.png}
\end{minipage}
\end{titlepage}
\frontmatter
\begin{center}
{\LARGE\bfseries\titulo}\\
\end{center}
\begin{center}
\autor\
\end{center}

\section*{Resumen}

\newpage
\begin{center}
{\LARGE\bfseries\tituloEng}\\
\end{center}
\begin{center}
\autor\
\end{center}

\section*{Abstract}

\newpage
\thispagestyle{empty}
\
\vspace{3cm}

\noindent\rule[-1ex]{\textwidth}{2pt}\\[4.5ex]

Yo, \textbf{\autor}, alumno de la titulación \textbf{\master} de la \textbf{\escuela\ de la \universidad}, autorizo la ubicación de la siguiente copia de mi Trabajo Fin de Máster (\textit{\titulo}) en la biblioteca del centro para que pueda ser consultada por las personas que lo deseen.

\bigskip
Además, este mismo trabajo es realizado bajo licencia \textbf{Creative Commons Attribution-ShareAlike 4.0} (\url{https://creativecommons.org/licenses/by-sa/4.0/}), dando permiso para copiarlo y redistribuirlo en cualquier medio o formato, también de adaptarlo de la forma que se quiera, pero todo esto siempre y cuando se reconozca la autoría y se distribuya con la misma licencia que el trabajo original. El documento en formato {\tt LaTeX} se puede encontrar en el siguiente repositorio de {\tt GitHub}: \url{https://github.com/germaaan/ProyectoDatosAbiertos}.

\vspace{4cm}

\noindent Fdo: \autor

\vspace{2cm}

\begin{flushright}
\ciudad, a \today
\end{flushright}

\newpage
\thispagestyle{empty}
\
\vspace{3cm}

\noindent\rule[-1ex]{\textwidth}{2pt}\\[4.5ex]

D. \textbf{\tutor}, profesor del \textbf{Departamento de Arquitectura y Tecnología de Computadores} de la \textbf{\universidad}.

\vspace{0.5cm}

\vspace{0.5cm}

\textbf{Informa:}

\vspace{0.5cm}

Que el presente trabajo, titulado \textit{\textbf{\titulo}}, ha sido realizado bajo su supervisión por \textbf{\autor}, y 
autoriza la defensa de dicho trabajo ante el tribunal que corresponda.

\vspace{0.5cm}

Y para que conste, expide y firma el presente informe en \ciudad\ a \today.

\vspace{1cm}

\textbf{El tutor:}

\vspace{5cm}

\noindent \textbf{\tutor}

\chapter*{Agradecimientos}
\thispagestyle{empty}

\vspace{1cm}

Agradecimientos.

\begingroup
\let\cleardoublepage\clearpage
  \tableofcontents
  %\listoffigures
  \listoftables
  %\lstlistoflistings
\endgroup

\newpage
\thispagestyle{empty}
\
\mainmatter
\chapter{Introducción}

Introducción.
\chapter{Objetivos}

Objetivos.
\chapter{Planificación inicial del Trabajo}

\section{Modelo de desarrollo}

Si bien no existe una metodología correcta para el desarrollo de ontologías, principalmente debido a que un punto muy influyente la percepción que tengamos nosotros mismos de la \textit{realidad} que queremos modelar, si que existen propuestas para llevar a cabo esta tarea.

\subsection{METHONTOLOGY}

Esta metodología divide el desarrollo de una ontología en las siguientes fases:

\begin{enumerate}
	\item \textbf{Construir el glosario de términos}: incluimos todos los términos relevantes del dominio (conceptos, instancias, atributos, relaciones) y sus descripciones en lenguaje natural.
	\item \textbf{Construir taxonomías de conceptos}: seleccionamos todos los conceptos que hemos definido en nuestro glosario y definimos con la taxonomía la jerarquía en estos conceptos.
	\item \textbf{Construir un diagrama de relaciones binarias}: consiste en establecer las relaciones entre los conceptos de la ontología y sus inversas.
	\item \textbf{Construir el diccionario de conceptos}: especificar qué propiedades describen cada concepto de la taxonomía, las relaciones del diagrama y las instancias de cada uno de los conceptos.
	\item \textbf{Describir en detalle las relaciones binarias}: especificando nombre, origen, destino, cardinalidad y relación inversa.
	\item \textbf{Describir en detalle los atributos de instancia}: especificando nombre, concepto al que pertenece, tipo de valor, rango de valores, cardinalidad.
	\item \textbf{Describir en detalle los atributos de clase}: especificando nombre, concepto donde se define, tipo de valor, rango de valores, cardinalidad.
	\item \textbf{Describir en detalle las constantes}: especificando nombre, tipo de valor, valor y unidad de medida (si es numérica).
	\item \textbf{Definir los axiomas formales}: especificando nombre, descripción de la restricción, expresión lógica a cumplir, conceptos a los que afecta, relaciones a las que afecta y variables que usa.
	\item \textbf{Definir reglas}: especificando nombre, descripción de la regla, expresión en forma si-entonces, conceptos a los que afecta, atributos a los que afecta, relaciones a las que afecta y variables que usa.
	\item \textbf{Describir instancias}: especificando nombre de la instancia, nombre del concepto al que pertenece y los valores de los atributos si se conocen.
\end{enumerate}

\subsection{On-To-Knowledge}

Esta metodología divide el desarrollo de una ontología en las siguientes fases:

\begin{enumerate}
	\item \textbf{Estudio de viabilidad}: se identifica el problema a resolver y sus posibles solucions.
	\item \textbf{Arranque}: se crea un documento de especificaciones de requisitos de la ontología: objetivo, dominio, alcance, aplicaciones que van a hacer uso de ella, fuentes de conocimiento, usuarios y escenarios, preguntas de competencia y ontologías existentes.
	\item \textbf{Refinamiento}: se realiza una clasificación de los términos de la ontología en vista de una posterior formalización con relaciones y axiomas.
	\item \textbf{Evaluación}: se comprueba la utilidad de ontología poniendo a prueba que cumple con los requisitos definidos en el documento de especificaciones.
	\item \textbf{Mantenimiento}: se actualiza la ontología en función de los cambios que se haga en las especificaciones durante el ciclo de vida de la propia ontología.
\end{enumerate}

\subsection{Protégé}

{\it Protégé} es un editor de ontologías de código abierto desarrollado por la {\it Universidad de Stanford} que permite desarrollar ontologías siguiendo una metología creada por los propios desarrolladores de {\it Protégé}. Se componente de las siguientes fases:

\begin{enumerate}
	\item \textbf{Determinar el dominio y el alcance de la ontología}: se establecen cuestiones como para qué se va a utilizar, qué preguntas ha de responder o a quién va dirigida.
	\item \textbf{Considerar la reutilización de ontologías existentes}: se busca reaprovechar recursos de ontologías existentes que se pueden considerar de utilidad para nuestro caso; además, de facilitar la interconexión con otras aplicaciones que hagan uso de otro ontologías.
	\item \textbf{Enumerar términos importantes para la ontología}: realizando una lista de conceptos que se quieren tratar con lo ontología.
	\item \textbf{Definir las clases y su jerarquía}: usando los conceptos de la lista anterior para dar forma a la taxonomía de la ontología.
	\item \textbf{Definir las propiedades de las clases}: para así darle capacidad de representación de información de las clases.
	\item \textbf{Definir las restricciones de las propiedades}: tipo de valores, cardinalidad, dominio, rango...
	\item \textbf{Crear instancias}: indicando los valores de las propiedades de cada una de las instancias de las clases.
\end{enumerate}
	
\subsection{Metodología seleccionada}

Debido a que en el problema que abarca este {\it Trabajo} no se parte de un sistema inexistente, si no que se va a trabajar sobre unos conceptos ya existentes, el uso de metodologías tan formales como {\it METHONTOLOGY} o {\it On-To-Knowledge} se hace algo de difícil de adaptación, es por eso que se han incluido como simples planteamientos teóricos.

\bigskip
En el caso de la ontología usada en {\it Protégé}, aunque existe el mismo problema, si es cierto que se puede adaptar a nuestro caso, por lo que es lo que se va a utilizar en el desarollo.

\section{Gestión de recursos}

\subsection{Personal}

El personal que se encargará del desarrollo de este proyecto es única y exclusivamente el autor del mismo, {\it Germán Martínez Maldonado}, encargándose de todas las fases del mismo: \textit{análisis}, \textit{diseño}, \textit{implementación} y \textit{evaluación de resultados}.

\subsection{Hardware}

El hardware utilizado para el desarrollo del proyecto se compone de dos máquinas diferentes: un ordenador portátil personal en el que se realizará el desarrollo y las pruebas haciendo uso de máquinas virtuales, y un servidor con visibilidad pública que será en el que se desplegará la versión final de proyecto. Las especificaciones de ambos sistemas son las siguientes.

\begin{table}[!ht]
	\centering
	\begin{tabular}{|p{.15\textwidth}|p{.35\textwidth}|p{.35\textwidth}|}
		\hline
		 &
		\textbf{Ordenador personal}&
		\textbf{Servidor de producción}
		\\ \hline
		CPU:&
		Intel Core i7-5500U CPU @ 3GHz&
		Intel Core2 Duo CPU E4400 @ 2GHz
		\\ \hline
		RAM:&
		8 GB&
		4 GB
		\\ \hline
		SO:&
		Ubuntu 16.10 yakkety&
		Ubuntu 16.04 xenial
		\\ \hline
		Kernel:&
		x86\_64 Linux 4.8.0-58-generic&
		x86\_64 Linux 4.4.0-83-generic
		\\ \hline
	\end{tabular}
	\caption{Características hardware utilizado}
	\label{caracteristicas-hardware}
\end{table}

\subsection{Software}

Todo el software que se ha utilizado para el desarrollo del proyecto es software de código abierto en un compromiso con la ciencia abierta y sostenible.

\begin{itemize}
	\item \textbf{Vagrant}\footnote{\url {https://www.vagrantup.com/}}: es una herramienta que nos permite crear fácilmente entornos virtuales en los que probar nuestro proyecto durante el desarollo. 
	\item \textbf{Ansible}\footnote{\url {https://www.ansible.com/}}: es una herramienta para automatizar la configuración de un sistema, instalando todo el software que necesitemos y permitiéndonos adaptarlo a nuestras necesidades. 
	\item \textbf{CKAN}\footnote{\url {https://ckan.org/}}: es una plataforma con la que podemos montar nuestro propio almacén de datos abiertos en el que publicar y visualizar datos en diferentes formatos.	
	\item \textbf{NGINX}\footnote{\url {https://nginx.org/en/}}: es un servidor web ligero, sencillo y ágil que usaremos para albergar los archivos de ontología y recursos producidos como fruto de este proyecto.
	\item \textbf{OpenLink Virtuoso Open-Source Edition}\footnote{\url {http://vos.openlinksw.com/owiki/wiki/VOS/}}: es un servidor ORDBMS abierto que permite el almacenamiento y la gestión de datos en formato RDF, además de proveer de un punto de acceso a un sistema de recuperación de datos mediante SPARQL.
\end{itemize}

\section{Planificación temporal}

El desarrollo de este {\tt Trabajo} se ha llevado a cabo entre principios de enero de 2017 y principios de julio de ese mismo año, habiéndose dividido en las siguientes fases:

\begin{enumerate}
	\item \textbf{Investigación inicial sobre la temática}: antes de empezar es necesario hacer un trabajo previo de investigación para adquirir todos los conocimientos que sean necesarios. metodologías para el desarrollo de ontologías, estándares de {\tt la Web Semántica}, endpoints {\tt SPARQL}...; aunque se irá viendo más en profundidad según se vaya avanzando en el {\tt Trabajo}, este trabajo previo es a lo que se dedicarán los primeros meses.
	\item \textbf{Documentación}: en este proyecto la documentación tiene un papel muy importante, ya que la parte de mayor complejidad corresponde al análisis y diseño del sistema que solo verán representadas en la misma; pero aunque la más importante, solo es una de las partes, así que con el fin de obtener una buena documentación, está se irá realizando simultáneamente con los avances que se vayan produciendo durante todo el desarrollo del proyecto.
	\item \textbf{Análisis y diseño del sistema}: como hemos dicho esta es la parte más importante del proyecto, ya que de ella depende el resto del desarrollo. Con el fin de estudiar todos los aspectos necesarios de la mejor forma posible se emplearán otro par de meses para esta tarea.
	\newpage
	\item \textbf{Desarrollo de los objetivos}: una vez que tenemos el diseño ya finalizado la implementación será bastante sencilla, ya que simplemente deberemos procesar los datos originales para que se adapten al diseño realizado. Con esto hecho solo nos quedará poner el marcha el servidor que albergará los datos y el servidor que proveerá el punto de acceso {\tt SPARQL}. Esta tarea no debería llevar más de un mes.
	\item \textbf{Análisis de resultados y conclusiones}: Con todo el trabajo ya finalizado, solo resta analizar los resultados para comprobar que cumplen los requisitos establecidos y, por fin, ya solo quedaría pasar a escribir las conclusiones y posibles trabajos futuros que se puedan derivar del trabajo realizado. El último mes se empleará en este fin.
\end{enumerate}
\chapter{Estándares de la Web Semántica}

\section{Ontologías}
Aunque en la introducción hemos explicado brevemente lo que era una ontología, lo que básicamente describíamos como una definición formal de los conceptos del dominio de interés con el que estamos tratando de indicar cuáles son sus tipos, propiedades y relaciones para conocer el estado completo de nuestro ``mundo". Pues para construir estas ontologías podremos usar dos lenguajes independientes, pero que que se suelen usar conjuntamente: RDF Schema y OWL.

\section{RDF Schema}
RDF Schema (o simplemente RDFS) es una extensión del RDF original que utiliza su misma estructura y que aparece para solucionar problemas básicos en la definición de ontologías como pueden ser la declaración de clases, las restricciones entre ellas y las restricciones de sus propiedades; ya que en la versión original de RDF solo eran posible definir tipos (lo que no es exactamente igual que una clase) y las propiedades de los mismos.

\bigskip
Gracias a que RDFs nos provee de los elementos básicos y comunes para la descripción de los datos de nuestro dominio, esto nos puede servir para reutilizar conceptos que son comunes en varios dominios distintos. Para conseguir esto RDFS nos permite crear esquemas sencillos usando clases y subclases, además de permitirnos definir propiedades, dominios de aplicación y rangos de valores.

\newpage
\begin{itemize}
	\item Las clases son conjuntos de recursos que tienen características comunes y una representación en el mundo real. Tenemos tres clases fundamentales en RDFs a partir de las que podemos definir nuestras propias clases, siempre posible crear además una jerarquía de clases.
	\begin{itemize}
		\item rdfs:Class es la clase que utilizaremos para definir otras clases.
		\item rdfs:Property es la clase a partir de la cual definiremos nuevas propiedades que nos permitan describir nuestros recursos.
		\item rdfs:Resource todas las cosas descritas en RDF son llamadas ``recursos" y son instancias de esta clase, es por eso que esta clase nos permite referenciar cualquier clases desde otra clase.
	\end{itemize}
	
	\item Las propiedades nos permitirán describir las características que tienen los distintos recursos, además de las propias relaciones que existen entre distintos recursos.
	\begin{itemize}
		\item rdf:type nos permite definir el tipo de un determinado recurso, además, un recursos puede ser una instancia de más de una clase.
		\item rdfs:subClassOf es la que nos permitirá definir una jerarquía de clases y subclases, además, una clase puede ser subclase de otras subclases.
		\item rdfs:subPropertyOf, al igual que podemos crear subclases, también podemos crear subpropiedades para establecer una jerarquía de las mismas.
	\end{itemize}
	
	\item Las restricciones son las que nos permiten definir las clases sobre la que pueden aplicarse determinadas propiedades y posibles valores:
	\begin{itemize}
		\item rdfs:domain definiría las clases sobre las que se puede aplicar una propiedad.
		\item rdfs:range definiría los valores que puede tener una propiedad.
	\end{itemize}
\end{itemize}

\section{OWL}

OWL (Ontology Web Language) es un lenguaje de especificación de ontologías extensión de RDF Schema, proporcionando un mayor significado y semántica a las ontologías, ya que nos permite definir restricciones sobre las propiedades (como por ejemplo, los valores que puede tomar una clase) y definir axiomas, sentencias que son siempre ciertas y que nos serán muy útiles a la hora de realizar restricciones (como por ejemplo, que un recurso no puede pertenecer a dos clases que sean disjuntas). Utilizando esta jerarquía de clases y propiedades que forman la ontología, OWL se basa en lógicas descriptivas para realizar razonamientos.

\newpage
Por otra parte, OWL se puede clasificar en tres tipos, según su nivel de expresividad:

\begin{itemize}
	\item OWL Lite. Es la versión más simple y nos provee únicamente de los elementos necesarios para definir una jerarquía de clases y propiedades con restricciones básicas como los tipos de valores, cardinalidades o propiedades de las características (inversa, transitiva, simétrica, funcional e inversa funcinal) entre otras.
	\item OWL DL. Posee mayor expresividad que la anterior, permitiendo además que las jerarquías sean razonadas automáticamente por inferencia que además permitan encontrar posibles inconsistencias mediante el uso de diferentes axiomas (como especificar clases disjuntas)
	\item OWL Full. Es la versión con un nivel mayor de expresividad, pero que está basadas en semánticas diferentes a OWL Lite y OWL DL. Por ejemplo, en OWL Full, una clase puede ser simultáneamente un colección de individuos y un individuo en sí mismo, algo que no está permitido en OWL DL.
\end{itemize}

Y después de hablar sobre los tipos de clasificación de OWL, tenemos que especificar sus componentes generales:

\begin{itemize}
	\item Clases e instancias: todas las clases (owl:Class) son a su vez subclases de una única superclase (owl:THing). Además las subclases e instancias se definen con RDFS mediante las propiedades rdfs:subClassOf y rdfs:type respectivamente. También una clase puede crearse como la intersección de varias clases (owl:intersectionOf), la unión de varias clases (owl:unionOf) o el complemento de otra clase (owl:complementOf).
	\item Propiedades: podemos definir propiedades literales de una clase con owl:DatatypeProperty, indicando además el tipo de dato de XML Schema del que es dicha propiedad; además, también podemos indicar las propiedades que relacionan distintas clases con owl:ObjectProperty.
	\item Características de las propiedades: OWL nos permite definir diferentes tipos de propiedades:
	\begin{itemize}
		\item Transitiva (owl:TransitiveProperty): si el sujeto A está relacionado con el sujeto B por medio de la propiedad P y por otro lado el sujeto B está relacionado con el sujeto C también por medio de la propiedad P, si la propiedad P está definida como una propiedad transitiva, se puede inferir que el sujeto A está relacionado con el sujeto C por medio de la propiedad P.
		\item Simétrica (owl:SymmetricProperty): si el sujeto A está relacionado con el sujeto B por medio de la propiedad P, si la propiedad P está definida como una propiedad simétrica, se puede inferir que el sujeto B está relacionado con el sujeto A por medio de la propiedad P.
		\item Funcional (owl:FunctionalProperty): si el sujeto A está relacionado con el sujeto B por medio de la propiedad P y por otro lado el sujeto B está relacionado con el sujeto C también por medio de la propiedad P, si la propiedad P está definida como una propiedad funcional, se puede inferir que el sujeto B y el sujeto C son el mismo sujeto.
		\item Inversa (owl:InverseOf): si la propiedad P1 está definida como la propiedad inversa de la propiedad P2, entonces si el sujeto A está relacionado con el sujeto B por medio de la propiedad P1, se puede inferir que el sujeto B está relacionado con el sujeto A por medio de la propiedad P2.
		\item Inversa funcional (owl:InverseFunctionalProperty): si la propiedad P está definida como propiedad inversa funcional, entonces si el sujeto A está relacionado con el sujeto B por medio de la propiedad P, se puede inferir que el sujeto A y el sujeto B son en realidad el mismo sujeto.
	\end{itemize}
	\item Restricciones sobre las propiedades: una vez definidas las clases y las propiedades solo nos queda ver como se definen las restricciones sobre las propiedades. Además del dominio y el rango que ya podíamos definir con RDFS, OWL nos permite hacer restricciones más descriptivas como indicar de qué forma las instancias de una clase pueden tomar como únicos valores las instancias de otra clase (owl:allValuesFrom, owl:someValuesFrom, owl:hasValue) o la cardinalidad de los mismos (owl:maxCardinality, owl:minCardinality, owl:cardinality).
	\item Axiomas: un axioma es una sentencia que es siempre cierta, OWL DL se basa en la lógica descriptiva para partiendo de la jerarquía y restricciones de clases poder inferir de forma automática razonamientos sobre los datos que tiene. En concreto, podemos definir desde controlar la integridad de nuestros datos mediante la definición de clases o propiedades disjuntas (owl:disjointWith, owl:AllDisjointClasses, owl:AllDisjointProperties); o indicar que nuestras clases y propiedades son equivalentes a clases y propiedades de otros vocabularios (owl:equivalentClass, owl:equivalentProperty). Esto último es lo que realmente le da utilidad a tener estructuras con un marco común, ya que es lo que permite la interoperabilidad, permiten que existan los datos enlazados.
\end{itemize}

\section{SPARQL}

\section{Razonamiento e inferencia}

``Se entiende por razonamiento a la facultad que permite resolver problemas, extraer conclusiones y aprender de manera consciente de los hechos, estableciendo conexiones causales y lógicas necesarias entre ellos."\footnote{https://es.wikipedia.org/wiki/Razonamiento}.

\bigskip
Hemos dicho que el principal objetivo es conseguir que los datos sean entendibles por las máquinas, que puedan razonar sobre ellos y sacar conclusiones; precisamente para eso nos basaremos en la lógica. La lógica estudia las condiciones bajo las cuales siguiendo una serie de pasos elementales se puede pasar de una serie de premisas a una conclusión, cuando estamos tratando con máquinas el problema que encontramos en su forma de procesar la lógica es que no dispone de los mecanismos adecuados para conocer cómo y en qué orden deben realizarse esos pasos elementales, algo que si es capaz de hacer la mente humana de forma natural. 

\bigskip
Esto es lo que intenta solucionar la inteligencia artificial, y en nuestro caso concreto es lo que queremos conseguir, que a través de las ontologías definidas se puedan realizar inferencias sobre la información proporcionada por la ontología descrita, que podamos obtener resultados producto del razonamiento que lleven a que la máquina saque sus propias conclusiones.

\bigskip
Por ejemplo, si tenemos varias clases que entre sus datos cuentan con la propiedad ``hombres" y la propiedad ``mujeres" en referencia al número de hombres o mujeres que se han matriculado en una determinada titulación, proceden de un mismo país o cualquier caso similar, nuestro razonamiento humano nos hace saber que los ``hombres" y las ``mujeres" son ``personas", por lo que si quisiéramos saber cuántas ``personas" hay matriculadas en una titulación, lo que estaríamos haciendo es sumar el número de ``hombres" y de ``mujeres"; sin embargo, este razonamiento no es obvio para una máquina, necesita una regla que le diga que los ``hombres" y las ``mujeres" son ``personas".

\bigskip
Aquí es donde vuelve a entrar en juego RDFS, ya que como hemos indicado antes, su vocabulario nos permite establecer jerarquías, por lo tanto podríamos especificar una jerarquía en la que tenemos una propiedad que es ``persona" y luego indicar que ``hombres" y ``mujeres" son subpropiedades de ``persona". Gracias a esta propiedad que acabamos de definir, si la máquina quiere contar número de ``personas", pero estas no están definidas explícitamente como tales, podrá inferir que lo que debe hacer es contar los ``hombres" y las ``mujeres".
\chapter{Análisis y Diseño}

\section{Conceptos del sistema}

Los conjuntos de datos con los que se ha decidido trabajar son los siguientes:

\begin{itemize}
	\item ``Demanda académica: procedimientos acceso": contiene información relativa al número total de solicitudes de matrícula demandadas en la universidad. Datos desde el curso 2012/2013 hasta el curso 2014/2015.
	\item ``Demanda académica: titulaciones": contiene información sobre la demanda de matrícula en relación con las plazas ofertadas en titulaciones oficiales de grado en la universidad. Datos desde el curso 2014/2015 hasta el curso 2014/2015.
	\item ``Matrículas: grado": contiene información relacionada con las matrículas de titulaciones de grado realizadas en la universidad, agrupándola por rama de conocimiento, titulación y sexo del estudiante. Datos desde el curso 2010/2011 hasta el curso 2014/2015.
	\item ``Matrículas: posgrado": contiene información relacionada con las matrículas de titulaciones de posgrado realizadas en la universidad, agrupándola por rama de conocimiento, titulación y sexo del estudiante. Datos desde el curso 2010/2011 hasta el curso 2014/2015.
	\item ``Número medio de créditos": contiene información relacionada con el número medio de créditos de los estudiantes de la universidad, agrupándola por rama de conocimiento, plan de estudios, número medio de créditos matriculados, número medio de créditos presentados y número medio de créditos superados. Datos desde el curso 2012/2013 hasta el curso 2013/2014.
	\item ``Oferta de titulaciones: doctorado": contiene información relativa a la oferta de titulaciones para estudios de doctorado, agrupándola por titulación, rama de conocimiento, centro y campus. Datos desde el curso 2013/2014 hasta el curso 2015/2016.
	\item ``Oferta de titulaciones: grado": contiene información relativa a la oferta de titulaciones para estudios de grado, agrupándola por titulación, rama de conocimiento, centro y campus. Datos desde el curso 2013/2014 hasta el curso 2015/2016.
	\item ``Oferta de titulaciones: másteres oficiales": contiene información relativa a la oferta de titulaciones para estudios de másteres oficiales, agrupándola por titulación, rama de conocimiento, centro y campus. Datos desde el curso 2013/2014 hasta el curso 2015/2016.
	\item ``Origen geográfico de estudiantes por país": contiene información relativa al origen geográfico de los estudiantes de la universidad, agrupándola por país de origen y sexo del estudiante. Datos desde el curso 2013/2014 hasta el curso 2015/2016.
	\item ``Origen geográfico de estudiantes por provincia": contiene información relativa al origen geográfico de los estudiantes de la universidad, agrupándola por provincia de origen y sexo del estudiante. Datos desde el curso 2013/2014 hasta el curso 2015/2016.
	\item ``Tasas académicas por titulaciones": contiene información relativa a la tasa de rendimiento, tasa de éxito, tasa de abandono inicial, tasa de eficiencia, tasa de graduación y tasa de abandono de los estudiantes según la titulación que estén estudiando en la universidad, agrupándola por titulación, tasa de rendimiento, tasa de éxito, tasa de abandono inicial, tasa de eficiencia, tasa de graduación y tasa de abandono. Datos desde el curso 2011/2012 hasta el curso 2015/2016.
\end{itemize}

El motivo de haber seleccionado solo estos 11 conjuntos de datos de entre los 40 totales, es que dado la gran cantidad de datos que podemos encontrar en el portal, podemos encontrarnos datos de todo tipo; podemos encontrar datos de fácil interpretación como son los datos relacionados con matrículas, pero también podemos encontranos por ejemplo datos de carácter económico que utiliza una serie de códigos que nos permiten hacer un uso tan natural de la información que contienen, por eso en una primera aproximación solo vamos a incluir datos de los que como personas podríamos sacar información válida a simple vista. Además, si no se tienen datos de varios años tampoco se considerará un conjunto de datos interesante ya que no se puede comprobar si esos datos representan un hecho anecdótico o una tendencia.

\bigskip

Aclarado este punto, lo siguiente es concretar que cada uno de estos conjuntos de datos se corresponderá con una clase en nuestra ontología.

\begin{itemize}
	\item ``Demanda académica: procedimientos acceso"  $\rightarrow$  ``DemandaAcademicaAcceso"
	\item ``Demanda académica: titulaciones"  $\rightarrow$  ``DemandaAcademicaTitulacion"
	\item ``Matrículas: grado"  $\rightarrow$  ``MatriculasGrado"
	\item ``Matrículas: posgrado"  $\rightarrow$  ``MatriculasPosgrado"
	\item ``Número medio de créditos"  $\rightarrow$  ``NumMedioCreditos"
	\item ``Oferta de titulaciones: doctorado"  $\rightarrow$  ``OfertaTitulacionDoctorado"
	\item ``Oferta de titulaciones: grado"  $\rightarrow$  ``OfertaTitulacionGrado"
	\item ``Oferta de titulaciones: másteres oficiales"  $\rightarrow$  ``OfertaTitulacionMaster"
	\item ``Origen geográfico de estudiantes por país"  $\rightarrow$  ``OrigenPais"
	\item ``Origen geográfico de estudiantes por provincia"  $\rightarrow$  ``OrigenProvincia"
	\item ``Tasas académicas por titulaciones"  $\rightarrow$  ``TasasAcademicasTitulacion"
\end{itemize}

\subsection{Definición de atributos, relaciones y restricciones}

Con las clases definidas, lo siguiente sería definir sus atributos, las relaciones y restricciones para tener la estructura completa de nuestro sistema. Para esto se usan básicamente dos tipos de propiedades diferentes definidas en OWL, ambas subclases de la clase rdf:Property definida en el estándar RDF:

\begin{itemize}
	\item Por un lado los atributos los definiremos con owl:DatatypeProperty, que es la propiedad que nos permite vincular valores del tipo que definamos con instancias de las clases que hemos definido.
	\item Por el otro lado las relaciones las definiremos con owl:ObjectProperty, que es la propiedad que nos permite vincular instancias de las distintas clases que hemos definido.
\end{itemize}

En lo que se refiere a los atributos, debemos tener en cuenta que hay atributos que se repiten en varias clases, por lo que tendremos que especificarlo en los dominios de dichas propiedades.

\begin{itemize}
	\item Campus: representa cada uno de los campus universitarios en los que está dividida la universidad.
	\item Centro: representa cada uno de los centros de la universidad, ya sean facultades, escuelas o centros adscritos.
	\item Cupo general: representa el número de matrículas realizadas por el grupo de estudiantes que pertenece al cupo general de estudiantes.
	\item Curso: representa el curso académico al que pertenecen los datos.
	\item Deportistas: representa el número de matrículas realizadas por el grupo de estudiantes que pertenece al cupo de deportistas de alto nivel o alto rendimiento.
	\item Discapacitados: representa el número de matrículas realizadas por el grupo de estudiantes que pertenece al cupo de personas con minusvalías reconocidas.
	\item Doctorado: representa las titulaciones de doctorado que se ofertan en la universidad.
	\item Estado: representa el estado de las solicitudes de matrículas demandadas en la universidad.
	\item Grado: representa las titulaciones de grado que se ofertan en la universidad.
	\item Hombres: representa el número de estudiantes del sexo masculino que estudian alguna titulación en la universidad.
	\item Máster: representa las titulaciones de másteres oficiales que se ofertan en la universidad.
	\item Mayores 25: representa el número de matrículas realizadas por el grupo de estudiantes que pertenece al cupo de personas mayores de 25 años.
	\item Mayores 40 y 45 : representa el número de matrículas realizadas por el grupo de estudiantes que pertenece al cupo de personas mayores de 40 y de 45 años.
	\item Mujeres: representa el número de estudiantes del sexo masculino que estudian alguna titulación en la universidad.
	\item Número medio de créditos matriculados:
	\item Número medio de créditos presentados:
	\item Número medio de créditos superados:
	\item País de origen: representa el número de estudiantes procedentes de un mismo país de origen que estudian alguna titulación en la universidad.
	\item Personas: representa el número de estudiantes que estudian alguna titulación en la universidad.
	\item Plan de estudios: representa el plan de estudios al que pertenece una titulación.
	\item Plazas ofertadas: representa el número de plazas ofertadas para una determinada titulación.
	\item Provincia: representa el número de estudiantes procedentes de una misma provincia del territorio español que estudian alguna titulación en la universidad.
	\item Rama: representa la rama de conocimiento a la que pertenece alguna titulación.
	\item Tasa de abandono: representa el porcentaje entre el número total de estudiantes de nuevo ingreso en una titulación que debieron obtener el título el año académico anterior y que no se han matriculado ni en ese año académico ni en el anterior.
	\item Tasa de abandono inicial: representa el porcentaje entre los estudiantes matriculados en una determinada titulación en un curso académico que no se matricularon en dicha titulación en los dos años siguientes y el número total de estudiantes que accedieron a esa misma titulación en ese mismo curso académico.
	\item Tasa de eficiencia: representa el porcentaje entre el número total de créditos de la titulación a los que deberían haberse matriculado el conjunto de estudiantes graduados en un año académico y el número total de créditos que se matricularon finalmente.
	\item Tasa de graduación: representa el porcentaje de estudiantes que finalizan una titulación en el tiempo previsto por el plan de estudios o en un año académico más y el número de estudiantes que entraron en esa misma titulación.
	\item Tasa de rendimiento: representa el porcentaje entre el número total de créditos superados (menos los créditos adaptados, convalidados y reconocidos) por los estudiantes de una titulación y el número total de créditos matriculados.
	\item Tipo de procedimiento: representa la forma de acceso por la que los estudiantes han realizado la solicitud de matrícula en la universidad.
	\item Titulación: representa las titulaciones de grado que se ofertan en la universidad.
	\item Titulados: representa el número de matrículas realizadas por el grupo de estudiantes que pertenece al cupo de titulados universitarios.
	\item Universidad: representa el conjunto de todos los campus de la universidad.
\end{itemize}

Con todos los atributos descritos ahora solo nos queda indicar cómo los vamos a referenciar en nuestra ontología y el tipo de dato que serán.

\begin{itemize}
	\item ``Campus" $\rightarrow$ ``campus". Tipo de dato cadena de caracteres.
	\item ``Centro" $\rightarrow$ ``centro". Tipo de dato cadena de caracteres.
	\item ``Cupo general" $\rightarrow$ ``cupoGral". Tipo de dato numérico entero no negativo.
	\item ``Curso" $\rightarrow$ ``curso". Tipo de dato cadena de caracteres.
	\item ``Deportistas" $\rightarrow$ ``deportistas". Tipo de dato numérico entero no negativo.
	\item ``Discapacitados" $\rightarrow$ ``discapacitados". Tipo de dato numérico entero no negativo.
	\item ``Doctorado" $\rightarrow$ ``doctorado". Tipo de dato cadena de caracteres.
	\item ``Estado" $\rightarrow$ ``estado". Tipo de dato cadena de caracteres.
	\item ``Grado" $\rightarrow$ ``grado". Tipo de dato cadena de caracteres.
	\item ``Hombres" $\rightarrow$ ``hombres". Tipo de dato numérico entero no negativo.
	\item ``Máster" $\rightarrow$ ``master". Tipo de dato cadena de caracteres.
	\item ``Mayores 25" $\rightarrow$ ``mayor25". Tipo de dato numérico entero no negativo.
	\item ``Mayores 40 y 45 " $\rightarrow$ ``mayor40". Tipo de dato numérico entero no negativo.
	\item ``Mujeres" $\rightarrow$ ``mujeres". Tipo de dato numérico entero no negativo.
	\item ``Número medio de créditos matriculados" $\rightarrow$ ``creditosMatriculados". Tipo de dato numérico real no negativo.
	\item ``Número medio de créditos presentados" $\rightarrow$ ``creditosPresentados". Tipo de dato numérico real no negativo.
	\item ``Número medio de créditos superados" $\rightarrow$ ``creditosSuperados". Tipo de dato numérico real no negativo.
	\item ``País de origen" $\rightarrow$ ``pais". Tipo de dato cadena de caracteres.
	\item ``Personas" $\rightarrow$ ``personas". Tipo de dato numérico entero no negativo.
	\item ``Plan de estudios" $\rightarrow$ ``planEstudios". Tipo de dato cadena de caracteres.
	\item ``Plazas ofertadas" $\rightarrow$ ``plazasOfertadas". Tipo de dato numérico entero no negativo.
	\item ``Provincia" $\rightarrow$ ``provincia". Tipo de dato cadena de caracteres.
	\item ``Rama de conocimiento" $\rightarrow$ ``ramaConocimiento". Tipo de dato cadena de caracteres.
	\item ``Tasa de abandono" $\rightarrow$ ``tasaAbandono". Tipo de dato numérico real no negativo.
	\item ``Tasa de abandono inicial" $\rightarrow$ ``tasaAbandonoInicial". Tipo de dato numérico real no negativo.
	\item ``Tasa de eficiencia" $\rightarrow$ ``tasaEficiencia". Tipo de dato numérico real no negativo.
	\item ``Tasa de graduación" $\rightarrow$ ``tasaGraduacion". Tipo de dato numérico real no negativo.
	\item ``Tasa de rendimiento" $\rightarrow$ ``tasaRendimiento". Tipo de dato numérico real no negativo.
	\item ``Tipo de procedimiento" $\rightarrow$ ``tipoProcedimiento". Tipo de dato cadena de caracteres.
	\item ``Titulación" $\rightarrow$ ``titulacion". Tipo de dato cadena de caracteres.
	\item ``Titulados" $\rightarrow$ ``titulados". Tipo de dato numérico entero no negativo
	\item ``Universidad" $\rightarrow$  ``universidad". Tipo de dato cadena de caracteres.
\end{itemize}

\section{Diseño de la ontología}

\subsection{Vocabularios usados}

Aunque en el diseño de la ontología todas las clases son del lenguaje definido propio, sí que se han usado otros vocabularios para definir los metadatos de la ontología (Dublin Core\footnote{\url{http://dublincore.org/documents/2012/06/14/dcmi-terms/}}, FOAF\footnote{\url{http://xmlns.com/foaf/spec/}}, Creative Commons\footnote{\url{https://creativecommons.org/ns}} y VANN\footnote{\url{http://vocab.org/vann/}}). \bigskip

Para la definición del propio vocabulario (cuyo prefijo de espacio de nombres preferido es \textbf{ugr}), solo se han utilizado los estándares definidos RDF, RDFS, OWL y XMLS\footnote{\url{https://www.w3.org/2001/XMLSchema}} (para los tipos de datos); sin embargo, para conseguir el objetivo de conseguir disponer de datos enlazados en las diferentes propiedades se referencian propiedades equivalentes dentro de otros vocabularios de nivel superior como DBpedia\footnote{\url{http://dbpedia.org/ontology/}} Wikidata\footnote{\url{https://www.wikidata.org/wiki/}} o Scheme.org\footnote{\url{http://schema.org/}}.

\subsection{Clases}

Las tablas 3.1 a 3.11 contienen la lista de propiedades de cada una de las clases, así como un ejemplo de una instancia de la clase en formato RDF/XML. 

\subsection{Propiedades}

Las tablas 3.12 a 3.41 contienen el dominio (las clases en las que se encuentran presentes), el rango (el tipo de dato), subpropiedad de (en caso de que lo sea), las propiedades equivalentes en otros vocabularios de nivel superior y su descripción dentro de la ontología en formato RDF/XML.
\bigskip

\begin{table}[!ht]
	\centering
	\begin{tabular}{|p{.17\textwidth}|p{.9\textwidth}|}
		\hline
		\multicolumn{2}{|l|}{Clase: \textbf{DemandaAcademicaAcceso}}
		\\ \hline
		Propiedades:&
		\begin{itemize}
			\item tipoProcedimiento
			\item estado
			\item hombres
			\item mujeres
			\item curso
		\end{itemize}
		\\ \hline
		Ejemplo:&
		\textless rdf:Description \newline
		\tab rdf:about=``/DemandaAcademicaAcceso/2012-2013\#1"\textgreater \newline
		\tab \textless rdf:type rdf:resource=``\#DemandaAcademicaAcceso"\ /\textgreater \newline
		\tab \textless ugr:tipoProcedimiento\textgreater \newline\tab\tab CONVOCATORIA ORDINARIA DE JUNIO: PRUEBA DE ACCESO A LA UNIVERSIDAD PARA ESTUDIANTES PROVENIENTES DE BACHILLERATO Y DE CICLOS FORMATIVOS DE GRADO SUPERIOR \newline\tab\textless /ugr:tipoProcedimiento\textgreater \newline
		\tab \textless ugr:estado\textgreater \newline\tab\tab PRESENTADOS FASE GENERAL\newline\tab\textless /ugr:estado\textgreater  \newline
		\tab \textless ugr:hombres rdf:datatype=``\&xsd;nonNegativeInteger"\textgreater \newline\tab\tab2264 \newline \tab \textless /ugr:hombres\textgreater \newline
		\tab \textless ugr:mujeres rdf:datatype=``\&xsd;nonNegativeInteger"\textgreater \newline\tab\tab2877 \newline \tab \textless /ugr:mujeres\textgreater  \newline
		\tab \textless ugr:curso\textgreater \newline\tab\tab2012/2013\newline\tab\textless /ugr:curso\textgreater  \newline
		\textless /rdf:Description\textgreater 
		\\ \hline
	\end{tabular}
	\caption{Clase DemandaAcademicaAcceso}
	\label{clase-demandaacademicaacceso}
\end{table}

\begin{table}[!ht]
	\centering
	\begin{tabular}{|p{.17\textwidth}|p{.9\textwidth}|}
		\hline
		\multicolumn{2}{|l|}{Clase: \textbf{DemandaAcademicaTitulacion}}
		\\ \hline
		Propiedades:&
		\begin{itemize}
			\item titulacion
			\item plazasOfertadas
			\item cupoGral
			\item mayor25
			\item mayor40
			\item titulados
			\item discapacitados
			\item deportistas
		\end{itemize}
		\\ \hline
		Ejemplo:&
		\textless rdf:Description \newline \tab rdf:about=``/DemandaAcademicaTitulacion/2014-2015\#1"\textgreater \newline
		\tab \textless rdf:type rdf:resource=``\#DemandaAcademicaTitulacion"\ /\textgreater \newline
		\tab \textless ugr:titulacion\textgreater \newline\tab\tab ADMINISTRACIÓN Y DIRECCIÓN DE EMPRESAS\newline\tab\textless /ugr:titulacion\textgreater \newline
		\tab \textless ugr:plazasOfertadas rdf:datatype=``\&xsd;nonNegativeInteger"\textgreater  \newline \tab \tab 281\newline\tab\textless /ugr:plazasOfertadas\textgreater \newline
		\tab \textless ugr:cupoGral rdf:datatype=``\&xsd;nonNegativeInteger"\textgreater \newline \tab \tab 271\newline\tab\textless /ugr:cupoGral\textgreater 
		\tab \newline \tab \textless ugr:mayor25 rdf:datatype=``\&xsd;nonNegativeInteger"\textgreater \newline \tab \tab 6\newline\tab\textless /ugr:mayor25\textgreater 
		\tab \newline \tab \textless ugr:mayor40 rdf:datatype=``\&xsd;nonNegativeInteger"\textgreater \newline \tab \tab 0\newline\tab\textless /ugr:mayor40\textgreater 
		\tab \newline \tab \textless ugr:titulados rdf:datatype=``\&xsd;nonNegativeInteger"\textgreater \newline \tab \tab 3\newline\tab\textless /ugr:titulados\textgreater 
		\tab \newline \tab 
		\textless ugr:discapacitados rdf:datatype=``\&xsd;nonNegativeInteger"\textgreater \newline \tab \tab 1\newline\tab\textless /ugr:discapacitados\textgreater 
		\tab \newline \tab \textless ugr:deportistas rdf:datatype=``\&xsd;nonNegativeInteger"\textgreater \newline \tab \tab 1\newline\tab\textless /ugr:deportistas\textgreater 
		\tab \newline \tab \textless ugr:curso\textgreater \newline\tab\tab2014/2015\newline\tab\textless /ugr:curso\textgreater \newline
		\textless /rdf:Description\textgreater 
		\\ \hline
	\end{tabular}
	\caption{Clase DemandaAcademicaTitulacion}
	\label{clase-demandaacademicatitulacion}
\end{table}

\begin{table}[!ht]
	\centering
	\begin{tabular}{|p{.17\textwidth}|p{.9\textwidth}|}
		\hline
		\multicolumn{2}{|l|}{Clase: \textbf{MatriculasGrado}}
		\\ \hline
		Propiedades:&
		\begin{itemize}
			\item ramaConocimiento
			\item titulacion
			\item hombres
			\item mujeres
			\item curso
		\end{itemize}
		\\ \hline
		Ejemplo:&
		\textless rdf:Description rdf:about=``/MatriculasGrado/2010-2011\#1"\textgreater \newline
		\tab \textless rdf:type rdf:resource=``\#MatriculasGrado"\ /\textgreater \newline
		\tab \textless ugr:ramaConocimiento\textgreater \newline \tab\tab ARTES Y HUMANIDADES\newline\tab\textless /ugr:ramaConocimiento\textgreater \newline
		\tab \textless ugr:titulacion\textgreater \newline\tab\tab GRADO EN BELLAS ARTES\newline\tab\textless /ugr:titulacion\textgreater \newline
		\tab \textless ugr:hombres rdf:datatype=``\&xsd;nonNegativeInteger"\textgreater \newline\tab\tab70\newline\tab\textless /ugr:hombres\textgreater 
		\tab \newline\tab\textless ugr:mujeres rdf:datatype=``\&xsd;nonNegativeInteger"\textgreater \newline\tab\tab159\newline\tab\textless /ugr:mujeres\textgreater 
		\tab \newline\tab\textless ugr:curso\textgreater \newline\tab\tab2010/2011\newline\tab\textless /ugr:curso\textgreater \newline
		\textless /rdf:Description\textgreater 
		\\ \hline
	\end{tabular}
	\caption{Clase MatriculasGrado}
	\label{clase-matriculasgrado}
\end{table}

\begin{table}[!ht]
	\centering
	\begin{tabular}{|p{.17\textwidth}|p{.9\textwidth}|}
		\hline
		\multicolumn{2}{|l|}{Clase: \textbf{MatriculasPosgrado}}
		\\ \hline
		Propiedades:&
		\begin{itemize}
			\item titulacion
			\item hombres
			\item mujeres
			\item curso
		\end{itemize}
		\\ \hline
		Ejemplo:&
		\textless rdf:Description \newline\tab  rdf:about=``/MatriculasPosgrado/2010-2011\#1"\textgreater \newline
		\tab \textless rdf:type rdf:resource=``\#MatriculasPosgrado"\ /\textgreater \newline
		\tab \textless ugr:titulacion\textgreater \newline\tab\tab MASTER ERASMUS MUNDUS EN EL COLOR EN LA INFORMATICA Y LA TECNOLOGIA DE LOS MEDIOS / ERASMUS MUNDUS IN COLOR IN INFORMATICS AND MEDIA TECHNOLOGY (CIMET) \newline\tab\textless /ugr:titulacion\textgreater \newline
		\tab \textless ugr:hombres rdf:datatype=``\&xsd;nonNegativeInteger"\textgreater \newline\tab\tab9\newline\tab\textless /ugr:hombres\textgreater 
		\tab \newline\tab\textless ugr:mujeres rdf:datatype=``\&xsd;nonNegativeInteger"\textgreater \newline\tab\tab9\newline\tab\textless /ugr:mujeres\textgreater 
		\tab \newline\tab\textless ugr:curso\textgreater \newline\tab\tab2010/2011\newline\tab\textless /ugr:curso\textgreater \newline
		\textless /rdf:Description\textgreater 
		\\ \hline
	\end{tabular}
	\caption{Clase MatriculasPosgrado}
	\label{clase-matriculasposgrado}
\end{table}

\begin{table}[!ht]
	\centering
	\begin{tabular}{|p{.17\textwidth}|p{.9\textwidth}|}
		\hline
		\multicolumn{2}{|l|}{Clase: \textbf{NumMedioCreditos}}
		\\ \hline
		Propiedades:&
		\begin{itemize}
			\item planEstudios
			\item ramaConocimiento
			\item creditosMatriculados
			\item creditosPresentados
			\item creditosSuperados
			\item curso
		\end{itemize}
		\\ \hline
		Ejemplo:&
		\textless rdf:Description \newline \tab rdf:about=``/NumMedioCreditos/2012-2013\#1"\textgreater \newline
		\tab \textless rdf:type rdf:resource=``\#NumMedioCreditos"\ /\textgreater 
		\tab \newline\tab\textless ugr:planEstudios\textgreater \newline\tab\tab PRIMER/SEGUNDO CICLO\newline\tab\textless /ugr:planEstudios\textgreater 
		\tab \newline\tab\textless ugr:ramaConocimiento\textgreater \newline\tab\tab ARTES Y HUMANIDADES\newline\tab\textless /ugr:ramaConocimiento\textgreater 
		\newline \tab \textless ugr:creditosMatriculados rdf:datatype=``\&xsd;decimal"\textgreater \newline\tab\tab 52.40\newline\tab\textless /ugr:creditosMatriculados\textgreater \newline
		\tab \textless ugr:creditosPresentados rdf:datatype=``\&xsd;decimal"\textgreater \newline\tab\tab 40.06\newline\tab\textless /ugr:creditosPresentados\textgreater \newline
		\tab \textless ugr:creditosSuperados rdf:datatype=``\&xsd;decimal"\textgreater \newline\tab\tab 35.08\newline\tab\textless /ugr:creditosSuperados\textgreater \newline
		\tab \textless ugr:curso\textgreater \newline\tab\tab 2012/2013\newline\tab\textless /ugr:curso\textgreater 
		\newline\textless /rdf:Description\textgreater 
		\\ \hline
	\end{tabular}
	\caption{Clase NumMedioCreditos}
	\label{clase-nummediocreditos}
\end{table}

\begin{table}[!ht]
	\centering
	\begin{tabular}{|p{.17\textwidth}|p{.9\textwidth}|}
		\hline
		\multicolumn{2}{|l|}{Clase: \textbf{OfertaTitulacionDoctorado}}
		\\ \hline
		Propiedades:&
		\begin{itemize}
			\item rama
			\item titulacion
			\item campus
			\item centro
			\item curso
		\end{itemize}
		\\ \hline
		Ejemplo:&
		\textless rdf:Description \newline\tab rdf:about=``/OfertaTitulacionDoctorado/2013-2014\#1"\textgreater \newline
		\tab \textless rdf:type rdf:resource=``\#OfertaTitulacionDoctorado"\ /\textgreater 
		\newline \tab \textless ugr:ramaConocimiento\textgreater \newline\tab\tab ARTES Y HUMANIDADES\newline\tab\textless /ugr:ramaConocimiento\textgreater 
		\newline\tab \textless ugr:titulacion\textgreater \newline\tab\tab PROGRAMA DE DOCTORADO EN BIOMEDICINA\newline\tab \textless/ugr:titulacion\textgreater 
		\newline\tab \textless ugr:campus\textgreater \newline\tab\tab ESCUELA DE DOCTORADO\newline\tab\textless /ugr:campus\textgreater 
		\newline\tab \textless ugr:centro\textgreater \newline\tab\tab UGR\newline\tab\textless /ugr:centro\textgreater 
		\newline\tab \textless ugr:curso\textgreater \newline\tab\tab 2013/2014\newline\tab\textless /ugr:curso\textgreater 
		\newline\textless /rdf:Description\textgreater 
		\\ \hline
	\end{tabular}
	\caption{Clase OfertaTitulacionDoctorado}
	\label{clase-ofertatitulaciondoctorado}
\end{table}

\begin{table}[!ht]
	\centering
	\begin{tabular}{|p{.17\textwidth}|p{.9\textwidth}|}
		\hline
		\multicolumn{2}{|l|}{Clase: \textbf{OfertaTitulacionGrado}}
		\\ \hline
		Propiedades:&
		\begin{itemize}
			\item ramaConocimiento
			\item titulacion
			\item campus
			\item centro
			\item curso
		\end{itemize}
		\\ \hline
		Ejemplo:&
		\textless rdf:Description \newline\tab rdf:about=``/OfertaTitulacionGrado/2013-2014\#1"\textgreater \newline
		\tab \textless rdf:type rdf:resource=``\#OfertaTitulacionGrado"\ /\textgreater 
		\newline \tab \textless ugr:ramaConocimiento\textgreater \newline\tab\tab CIENCIAS\newline\tab\textless /ugr:ramaConocimiento\textgreater 
		\newline\tab \textless ugr:titulacion\textgreater \newline\tab\tab BIOLOGIA\newline\tab \textless/ugr:titulacion\textgreater 
		\newline\tab \textless ugr:campus\textgreater \newline\tab\tab FUENTENUEVA\newline\tab\textless /ugr:campus\textgreater 
		\newline\tab \textless ugr:centro\textgreater \newline\tab\tab FACULTAD DE CIENCIAS\newline\tab\textless /ugr:centro\textgreater 
		\newline\tab \textless ugr:curso\textgreater \newline\tab\tab 2013/2014\newline\tab\textless /ugr:curso\textgreater 
		\newline\textless /rdf:Description\textgreater 
		\\ \hline
	\end{tabular}
	\caption{Clase OfertaTitulacionGrado}
	\label{clase-ofertatitulaciongrado}
\end{table}

\begin{table}[!ht]
	\centering
	\begin{tabular}{|p{.17\textwidth}|p{.9\textwidth}|}
		\hline
		\multicolumn{2}{|l|}{Clase: \textbf{OfertaTitulacionMaster}}
		\\ \hline
		Propiedades:&
		\begin{itemize}
			\item ramaConocimiento
			\item titulacion
			\item campus
			\item centro
			\item curso
		\end{itemize}
		\\ \hline
		Ejemplo:&
		\textless rdf:Description \newline\tab rdf:about=``/OfertaTitulacionMaster/2013-2014\#1"\textgreater \newline
		\tab \textless rdf:type rdf:resource=``\#OfertaTitulacionMaster"\ /\textgreater 
		\newline \tab \textless ugr:ramaConocimiento\textgreater \newline\tab\tab  ARTES Y HUMANIDADES\newline\tab\textless /ugr:ramaConocimiento\textgreater 
		\newline\tab \textless ugr:titulacion\textgreater \newline\tab\tab MASTER UNIVERSITARIO EN ARQUEOLOGIA (M71.56.1)\newline\tab \textless/ugr:titulacion\textgreater 
		\newline\tab \textless ugr:campus\textgreater \newline\tab\tab CARTUJA\newline\tab\textless /ugr:campus\textgreater 
		\newline\tab \textless ugr:centro\textgreater \newline\tab\tab FACULTAD DE FILOSOFÍA Y LETRAS\newline\tab\textless /ugr:centro\textgreater 
		\newline\tab \textless ugr:curso\textgreater \newline\tab\tab 2013/2014\newline\tab\textless /ugr:curso\textgreater 
		\newline\textless /rdf:Description\textgreater 
		\\ \hline
	\end{tabular}
	\caption{Clase OfertaTitulacionMaster}
	\label{clase-ofertatitulacionmaster}
\end{table}

\begin{table}[!ht]
	\centering
	\begin{tabular}{|p{.17\textwidth}|p{.9\textwidth}|}
		\hline
		\multicolumn{2}{|l|}{Clase: \textbf{OrigenPais}}
		\\ \hline
		Propiedades:&
		\begin{itemize}
			\item pais
			\item hombres
			\item mujeres
			\item curso
		\end{itemize}
		\\ \hline
		Ejemplo:&
		\textless rdf:Description \newline\tab rdf:about=``/OrigenPais/2013-2014\#1"\textgreater 
		\tab \newline\tab \textless rdf:type rdf:resource=``\#OrigenPais"\ /\textgreater 
		\newline \tab \textless ugr:pais\textgreater \newline\tab\tab ALBANIA\newline\tab\textless /ugr:pais\textgreater 
		\newline\tab \textless ugr:hombres rdf:datatype=``\&xsd;nonNegativeInteger"\textgreater \newline\tab\tab 3\newline\tab\textless /ugr:hombres\textgreater 
		\newline\tab \textless ugr:mujeres rdf:datatype=``\&xsd;nonNegativeInteger"\textgreater \newline\tab\tab 2\newline\tab\textless /ugr:mujeres\textgreater 
		\newline\tab \textless ugr:curso\textgreater \newline\tab\tab 2013/2014\newline\tab\textless /ugr:curso\textgreater 
		\newline\textless /rdf:Description\textgreater 
		\\ \hline
	\end{tabular}
	\caption{Clase OrigenPais}
	\label{clase-origenpais}
\end{table}

\begin{table}[!ht]
	\centering
	\begin{tabular}{|p{.17\textwidth}|p{.9\textwidth}|}
		\hline
		\multicolumn{2}{|l|}{Clase: \textbf{OrigenProvincia}}
		\\ \hline
		Propiedades:&
		\begin{itemize}
			\item provincia
			\item hombres
			\item mujeres
			\item curso
		\end{itemize}
		\\ \hline
		Ejemplo:&
		\textless rdf:Description \newline\tab rdf:about=``/OrigenProvincia/2013-2014\#1"\textgreater 
		\tab \newline\tab \textless rdf:type rdf:resource=``\#OrigenProvincia"\ /\textgreater 
		\newline \tab \textless ugr:pais\textgreater \newline\tab\tab ALAVA\newline\tab\textless /ugr:pais\textgreater 
		\newline\tab \textless ugr:hombres rdf:datatype=``\&xsd;nonNegativeInteger"\textgreater \newline\tab\tab 24\newline\tab\textless /ugr:hombres\textgreater 
		\newline\tab \textless ugr:mujeres rdf:datatype=``\&xsd;nonNegativeInteger"\textgreater \newline\tab\tab 47\newline\tab\textless /ugr:mujeres\textgreater 
		\newline\tab \textless ugr:curso\textgreater \newline\tab\tab 2013/2014\newline\tab\textless /ugr:curso\textgreater 
		\newline\textless /rdf:Description\textgreater 
		\\ \hline
	\end{tabular}
	\caption{Clase OrigenProvincia}
	\label{clase-origenprovincia}
\end{table}

\begin{table}[!ht]
	\centering
	\begin{tabular}{|p{.17\textwidth}|p{.9\textwidth}|}
		\hline
		\multicolumn{2}{|l|}{Clase: \textbf{TasasAcademicasTitulacion}}
		\\ \hline
		Propiedades:&
		\begin{itemize}
			\item titulacion
			\item tasaRendimiento
			\item tasaExito
			\item tasaAbandonoInicial
			\item tasaEficiencia
			\item tasaGraduacion
			\item tasaAbandono
		\end{itemize}
		\\ \hline
		Ejemplo:&
		\textless rdf:Description\newline\tab rdf:about=``/TasasAcademicasTitulacion/2015-2016\#1"\textgreater 
		\newline\tab \textless rdf:type rdf:resource=``\#TasasAcademicasTitulacion"\ /\textgreater 
		\newline\tab \textless ugr:titulacion\textgreater \newline\tab\tab GRADUADO EN BIOLOGIA\newline\tab\textless /ugr:titulacion\textgreater 
		\newline\tab \textless ugr:tasaRendimiento rdf:datatype=``\&xsd;decimal"\textgreater \newline\tab\tab 74.02\newline\tab\textless /ugr:tasaRendimiento\textgreater 
		\newline\tab \textless ugr:tasaExito rdf:datatype=``\&xsd;decimal"\textgreater \newline\tab\tab 83.12\newline\tab\textless /ugr:tasaExito\textgreater 
		\newline\tab \textless ugr:tasaAbandonoInicial rdf:datatype=``\&xsd;decimal"\textgreater \newline\tab\tab 12.27\newline\tab\textless /ugr:tasaAbandonoInicial\textgreater 
		\newline\tab \textless ugr:tasaEficiencia rdf:datatype=``\&xsd;decimal"\textgreater \newline\tab\tab 98\newline\tab\textless /ugr:tasaEficiencia\textgreater 
		\newline\tab \textless ugr:tasaGraduacion rdf:datatype=``\&xsd;decimal"\textgreater \newline\tab\tab 32.73\newline\tab\textless /ugr:tasaGraduacion\textgreater 
		\newline\tab \textless ugr:tasaAbandono rdf:datatype=``\&xsd;decimal"\textgreater \newline\tab\tab30\newline\tab\textless /ugr:tasaAbandono\textgreater 
		\newline\tab \textless ugr:curso\textgreater \newline\tab\tab 2015/2016\newline\tab\textless /ugr:curso\textgreater 
		\newline\textless /rdf:Description\textgreater 
		\\ \hline
	\end{tabular}
	\caption{Clase TasasAcademicasTitulacion}
	\label{clase-tasasacademicastitulacion}
\end{table}

\begin{table}[!ht]
	\centering
	\begin{tabular}{|p{.25\textwidth}|p{.9\textwidth}|}
		\hline
		\multicolumn{2}{|l|}{Propiedad: \textbf{campus}}
		\\ \hline
		Dominio:&
		\begin{itemize}
			\item OfertaTitulacionDoctorado
			\item OfertaTitulacionGrado
			\item OfertaTitulacionMaster
		\end{itemize}
		\\ \hline
		Rango:&
		http://www.w3.org/2001/XMLSchema\#string
		\\ \hline
		Subpropiedad de:&
		universidad
		\\ \hline
		Propiedades \newline equivalentes:&
		\begin{itemize}
			\item \url{http://dbpedia.org/ontology/campus}
			\item \url{https://www.wikidata.org/wiki/Q209465}
		\end{itemize}
		\\ \hline
		Descripción:&
		\textless owl:DatatypeProperty rdf:about=``\#campus"\textgreater\newline 
		\tab\textless rdfs:label xml:lang=``es"\textgreater\newline
		\tab\tab Campus\newline
		\tab\textless /rdfs:label\textgreater\newline
		\tab\textless rdfs:subPropertyOf\newline
		\tab\tab rdf:resource=``\#universidad"\ /\textgreater\newline
		\tab\textless rdfs:range\newline
		\tab\tab rdf:resource=``\&xsd;string"\ /\textgreater\newline
		\tab\textless rdfs:domain\newline
		\tab\tab rdf:resource=``\#OfertaTitulacionDoctorado"\ /\textgreater\newline
		\tab\textless rdfs:domain\newline
		\tab\tab rdf:resource=``\#OfertaTitulacionGrado"\ /\textgreater\newline
		\tab\textless rdfs:domain\newline
		\tab\tab rdf:resource=``\#OfertaTitulacionMaster"\ /\textgreater\newline
		\tab\textless owl:equivalentProperty\newline
		\tab\tab rdf:resource=``http://dbpedia.org/ontology/campus"\ /\textgreater\newline
		\tab\textless owl:equivalentProperty\newline
		\tab\tab rdf:resource=``https://www.wikidata.org/wiki/Q209465"\ /\textgreater\newline
		\textless /owl:DatatypeProperty\textgreater
		\\ \hline
	\end{tabular}
	\caption{Propiedad campus}
	\label{propiedad-campus}
\end{table}

\begin{table}[!ht]
	\centering
	\begin{tabular}{|p{.25\textwidth}|p{.9\textwidth}|}
		\hline
		\multicolumn{2}{|l|}{Propiedad: \textbf{centro}}
		\\ \hline
		Dominio:&
		\begin{itemize}
			\item OfertaTitulacionDoctorado
			\item OfertaTitulacionGrado
			\item OfertaTitulacionMaster
		\end{itemize}
		\\ \hline
		Rango:&
		http://www.w3.org/2001/XMLSchema\#string
		\\ \hline
		Subpropiedad de:&
		campus
		\\ \hline
		Descripción:&
		\textless owl:DatatypeProperty rdf:about=``\#centro"\textgreater\newline 
		\tab\textless rdfs:label xml:lang=``es"\textgreater\newline
		\tab\tab Centro\newline
		\tab\textless /rdfs:label\textgreater\newline
		\tab\textless rdfs:subPropertyOf\newline
		\tab\tab rdf:resource=``\#campus"\ /\textgreater\newline
		\tab\textless rdfs:range\newline
		\tab\tab rdf:resource=``\&xsd;string"\ /\textgreater\newline
		\tab\textless rdfs:domain\newline
		\tab\tab rdf:resource=``\#OfertaTitulacionDoctorado"\ /\textgreater\newline
		\tab\textless rdfs:domain\newline
		\tab\tab rdf:resource=``\#OfertaTitulacionGrado"\ /\textgreater\newline
		\tab\textless rdfs:domain\newline
		\tab\tab rdf:resource=``\#OfertaTitulacionMaster"\ /\textgreater\newline
		\textless /owl:DatatypeProperty\textgreater
		\\ \hline
	\end{tabular}
	\caption{Propiedad centro}
	\label{propiedad-centro}
\end{table}

\begin{table}[!ht]
	\centering
	\begin{tabular}{|p{.25\textwidth}|p{.9\textwidth}|}
		\hline
		\multicolumn{2}{|l|}{Propiedad: \textbf{cupoGral}}
		\\ \hline
		Dominio:&
		\begin{itemize}
			\item DemandaAcademicaTitulacion
		\end{itemize}
		\\ \hline
		Rango:&
		http://www.w3.org/2001/XMLSchema\#decimal
		\\ \hline
		Descripción:&
		\textless owl:DatatypeProperty rdf:about=``\#cupoGral"\textgreater\newline 
		\tab\textless rdfs:label xml:lang=``es"\textgreater\newline
		\tab\tab Cupo general\newline
		\tab\textless /rdfs:label\textgreater\newline
		\tab\textless rdfs:range\newline
		\tab\tab rdf:resource=``\&xsd;decimal"\ /\textgreater\newline
		\tab\textless rdfs:domain\newline
		\tab\tab rdf:resource=``\#DemandaAcademicaTitulacion"\ /\textgreater\newline
		\textless /owl:DatatypeProperty\textgreater
		\\ \hline
	\end{tabular}
	\caption{Propiedad cupoGral}
	\label{propiedad-cupogral}
\end{table}

\begin{table}[!ht]
	\centering
	\begin{tabular}{|p{.25\textwidth}|p{.9\textwidth}|}
		\hline
		\multicolumn{2}{|l|}{Propiedad: \textbf{curso}}
		\\ \hline
		Dominio:&
		Todas las clases de ``ugr"
		\\ \hline
		Rango:&
		http://www.w3.org/2001/XMLSchema\#string
		\\ \hline
		Propiedades \newline equivalentes:&
		\begin{itemize}
			\item \url{http://purl.org/dc/terms/coverage}
		\end{itemize}
		\\ \hline
		Descripción:&
		\textless owl:DatatypeProperty rdf:about=``\#curso"\textgreater\newline 
		\tab\textless rdfs:label xml:lang=``es"\textgreater\newline
		\tab\tab Curso\newline
		\tab\textless /rdfs:label\textgreater\newline
		\tab\textless rdfs:range\newline
		\tab\tab rdf:resource=``\&xsd;string"\ /\textgreater\newline
		\tab\textless rdfs:domain\newline
		\tab\tab rdf:resource=``\#DemandaAcademicaAcceso"\ /\textgreater\newline
		\tab\textless rdfs:domain\newline
		\tab\tab rdf:resource=``\#DemandaAcademicaTitulacion"\ /\textgreater\newline
		\tab\textless rdfs:domain\newline
		\tab\tab rdf:resource=``\#MatriculasGrado"\ /\textgreater\newline
		\tab\textless rdfs:domain\newline
		\tab\tab rdf:resource=``\#MatriculasPosgrado"\ /\textgreater\newline
		\tab\textless rdfs:domain\newline
		\tab\tab rdf:resource=``\#NumMedioCreditos"\ /\textgreater\newline
		\tab\textless rdfs:domain\newline
		\tab\tab rdf:resource=``\#OfertaTitulacionDoctorado"\ /\textgreater\newline
		\tab\textless rdfs:domain\newline
		\tab\tab rdf:resource=``\#OfertaTitulacionGrado"\ /\textgreater\newline
		\tab\textless rdfs:domain\newline
		\tab\tab rdf:resource=``\#OfertaTitulacionMaster"\ /\textgreater\newline
		\tab\textless rdfs:domain\newline
		\tab\tab rdf:resource=``\#OrigenPais"\ /\textgreater\newline
		\tab\textless rdfs:domain\newline
		\tab\tab rdf:resource=``\#OrigenProvincia"\ /\textgreater\newline
		\tab\textless rdfs:domain\newline
		\tab\tab rdf:resource=``\#TasasAcademicasTitulacion"\ /\textgreater\newline
		\tab\textless owl:equivalentProperty\newline
		\tab\tab rdf:resource=``http://purl.org/dc/terms/coverage"\  /\textgreater\newline
		\textless /owl:DatatypeProperty\textgreater
		\\ \hline
	\end{tabular}
	\caption{Propiedad curso}
	\label{propiedad-curso}
\end{table}

\begin{table}[!ht]
	\centering
	\begin{tabular}{|p{.25\textwidth}|p{.9\textwidth}|}
		\hline
		\multicolumn{2}{|l|}{Propiedad: \textbf{deportistas}}
		\\ \hline
		Dominio:&
		\begin{itemize}
			\item DemandaAcademicaTitulacion
		\end{itemize}
		\\ \hline
		Rango:&
		http://www.w3.org/2001/XMLSchema\#nonNegativeInteger
		\\ \hline
		Propiedades \newline equivalentes:&
		\begin{itemize}
			\item \url{http://dbpedia.org/ontology/Athlete}
			\item \url{http://schema.org/athlete}
			\item \url{https://www.wikidata.org/wiki/Q2066131}
		\end{itemize}
		\\ \hline
		Descripción:&
		\textless owl:DatatypeProperty rdf:about=``\#deportistas"\textgreater\newline 
		\tab\textless rdfs:label xml:lang=``es"\textgreater\newline
		\tab\tab Deportistas\newline
		\tab\textless /rdfs:label\textgreater\newline
		\tab\textless rdfs:range\newline
		\tab\tab rdf:resource=``\&xsd;nonNegativeInteger"\ /\textgreater\newline
		\tab\textless rdfs:domain\newline
		\tab\tab rdf:resource=``\#DemandaAcademicaTitulacion"\ /\textgreater\newline
		\tab\textless owl:equivalentProperty\newline
		\tab\tab rdf:resource=``http://dbpedia.org/ontology/Athlete"\  /\textgreater\newline
		\tab\textless owl:equivalentProperty\newline
		\tab\tab rdf:resource=``http://schema.org/athlete"\  /\textgreater\newline
		\tab\textless owl:equivalentProperty\newline
		\tab\tab rdf:resource=``https://www.wikidata.org/wiki/Q2066131"\  /\textgreater\newline
		\textless /owl:DatatypeProperty\textgreater
		\\ \hline
	\end{tabular}
	\caption{Propiedad deportistas}
	\label{propiedad-deportistas}
\end{table}


\begin{table}[!ht]
	\centering
	\begin{tabular}{|p{.25\textwidth}|p{.9\textwidth}|}
		\hline
		\multicolumn{2}{|l|}{Propiedad: \textbf{discapacitados}}
		\\ \hline
		Dominio:&
		\begin{itemize}
			\item DemandaAcademicaTitulacion
		\end{itemize}
		\\ \hline
		Rango:&
		http://www.w3.org/2001/XMLSchema\#nonNegativeInteger
		\\ \hline
		Propiedades \newline equivalentes:&
		\begin{itemize}
			\item \url{https://www.wikidata.org/wiki/Q15978181}
		\end{itemize}
		\\ \hline
		Descripción:&
		\textless owl:DatatypeProperty rdf:about=``\#discapacitados"\textgreater\newline 
		\tab\textless rdfs:label xml:lang=``es"\textgreater\newline
		\tab\tab Discapacitados\newline
		\tab\textless /rdfs:label\textgreater\newline
		\tab\textless rdfs:range\newline
		\tab\tab rdf:resource=``\&xsd;nonNegativeInteger"\ /\textgreater\newline
		\tab\textless rdfs:domain\newline
		\tab\tab rdf:resource=``\#DemandaAcademicaTitulacion"\ /\textgreater\newline
		\tab\textless owl:equivalentProperty\newline
		\tab\tab rdf:resource=``https://www.wikidata.org/wiki/Q15978181"\  /\textgreater\newline
		\textless /owl:DatatypeProperty\textgreater
		\\ \hline
	\end{tabular}
	\caption{Propiedad discapacitados}
	\label{propiedad-discapacitados}
\end{table}

\begin{table}[!ht]
	\centering
	\begin{tabular}{|p{.25\textwidth}|p{.9\textwidth}|}
		\hline
		\multicolumn{2}{|l|}{Propiedad: \textbf{doctorado}}
		\\ \hline
		Dominio:&
		\begin{itemize}
			\item OfertaTitulacionDoctorado
		\end{itemize}
		\\ \hline
		Rango:&
		http://www.w3.org/2001/XMLSchema\#string
		\\ \hline
		Subpropiedad de:&
		titulacion
		\\ \hline
		Propiedades \newline equivalentes:&
		\begin{itemize}
			\item \url{http://dbpedia.org/page/Doctorate}
			\item \url{https://www.wikidata.org/wiki/Q849697}
		\end{itemize}
		\\ \hline
		Descripción:&
		\textless owl:DatatypeProperty rdf:about=``\#doctorado"\textgreater\newline 
		\tab\textless rdfs:label xml:lang=``es"\textgreater\newline
		\tab\tab Doctorado\newline
		\tab\textless /rdfs:label\textgreater\newline
		\tab\textless rdfs:subPropertyOf\newline
		\tab\tab rdf:resource=``\#titulacion"\ /\textgreater\newline
		\tab\textless rdfs:range\newline
		\tab\tab rdf:resource=``\&xsd;string"\ /\textgreater\newline
		\tab\textless rdfs:domain\newline
		\tab\tab rdf:resource=``\#OfertaTitulacionDoctorado"\ /\textgreater\newline
		\tab\textless rdfs:domain\newline
		\tab\textless owl:equivalentProperty\newline
		\tab\tab rdf:resource=``http://dbpedia.org/page/Doctorate"\  /\textgreater\newline
		\tab\textless owl:equivalentProperty\newline
		\tab\tab rdf:resource=``https://www.wikidata.org/wiki/Q849697"\  /\textgreater\newline
		\textless /owl:DatatypeProperty\textgreater
		\\ \hline
	\end{tabular}
	\caption{Propiedad doctorado}
	\label{propiedad-doctorado}
\end{table}

\begin{table}[!ht]
	\centering
	\begin{tabular}{|p{.25\textwidth}|p{.9\textwidth}|}
		\hline
		\multicolumn{2}{|l|}{Propiedad: \textbf{estado}}
		\\ \hline
		Dominio:&
		\begin{itemize}
			\item DemandaAcademicaAcceso
		\end{itemize}
		\\ \hline
		Rango:&
		http://www.w3.org/2001/XMLSchema\#string
		\\ \hline
		Descripción:&
		\textless owl:DatatypeProperty rdf:about=``\#estado"\textgreater\newline 
		\tab\textless rdfs:label xml:lang=``es"\textgreater\newline
		\tab\tab Estado\newline
		\tab\textless /rdfs:label\textgreater\newline
		\tab\textless rdfs:range\newline
		\tab\tab rdf:resource=``\&xsd;string"\ /\textgreater\newline
		\tab\textless rdfs:domain\newline
		\tab\tab rdf:resource=``\#DemandaAcademicaAcceso"\ /\textgreater\newline
		\textless /owl:DatatypeProperty\textgreater
		\\ \hline
	\end{tabular}
	\caption{Propiedad estado}
	\label{propiedad-estado}
\end{table}

\begin{table}[!ht]
	\centering
	\begin{tabular}{|p{.25\textwidth}|p{.9\textwidth}|}
		\hline
		\multicolumn{2}{|l|}{Propiedad: \textbf{grado}}
		\\ \hline
		Dominio:&
		\begin{itemize}
			\item DemandaAcademicaTitulacion
			\item MatriculasGrado
			\item OfertaTitulacionGrado
			\item TasasAcademicasTitulacion
		\end{itemize}
		\\ \hline
		Rango:&
		http://www.w3.org/2001/XMLSchema\#string
		\\ \hline
		Subpropiedad de:&
		titulacion
		\\ \hline
		Propiedades \newline equivalentes:&
		\begin{itemize}
			\item \url{http://dbpedia.org/page/Bachelor\%27s_degree}
			\item \url{https://www.wikidata.org/wiki/Q6008527}
		\end{itemize}
		\\ \hline
		Descripción:&
		\textless owl:DatatypeProperty rdf:about=``\#grado"\textgreater\newline 
		\tab\textless rdfs:label xml:lang=``es"\textgreater\newline
		\tab\tab Grado\newline
		\tab\textless /rdfs:label\textgreater\newline
		\tab\textless rdfs:subPropertyOf\newline
		\tab\tab rdf:resource=``\#titulacion"\ /\textgreater\newline
		\tab\textless rdfs:range\newline
		\tab\tab rdf:resource=``\&xsd;string"\ /\textgreater\newline
		\tab\textless rdfs:domain\newline
		\tab\tab rdf:resource=``\#DemandaAcademicaTitulacion"\ /\textgreater\newline
		\tab\textless rdfs:domain\newline
		\tab\tab rdf:resource=``\#MatriculasGrado"\ /\textgreater\newline
		\tab\textless rdfs:domain\newline
		\tab\tab rdf:resource=``\#OfertaTitulacionGrado"\ /\textgreater\newline
		\tab\textless rdfs:domain\newline
		\tab\tab rdf:resource=``\#TasasAcademicasTitulacion"\ /\textgreater\newline
		\tab\textless rdfs:domain\newline
		\tab\textless owl:equivalentProperty\newline
		\tab\tab rdf:resource=``http://dbpedia.org/page/Bachelor\%27s\_degree"\  /\textgreater\newline
		\tab\textless owl:equivalentProperty\newline
		\tab\tab rdf:resource=``https://www.wikidata.org/wiki/Q6008527"\  /\textgreater\newline
		\textless /owl:DatatypeProperty\textgreater
		\\ \hline
	\end{tabular}
	\caption{Propiedad grado}
	\label{propiedad-grado}
\end{table}

\begin{table}[!ht]
	\centering
	\begin{tabular}{|p{.25\textwidth}|p{.9\textwidth}|}
		\hline
		\multicolumn{2}{|l|}{Propiedad: \textbf{hombres}}
		\\ \hline
		Dominio:&
		\begin{itemize}
			\item DemandaAcademicaAcceso
			\item MatriculasGrado
			\item MatriculasPosgrado
			\item OrigenPais
			\item OrigenProvincia
		\end{itemize}
		\\ \hline
		Rango:&
		http://www.w3.org/2001/XMLSchema\#nonNegativeInteger
		\\ \hline
		Subpropiedad de:&
		personas
		\\ \hline
		Propiedades \newline equivalentes:&
		\begin{itemize}
			\item \url{http://dbpedia.org/page/Man}
			\item \url{https://schema.org/Male}
			\item \url{https://www.wikidata.org/wiki/Q8441}
		\end{itemize}
		\\ \hline
		Descripción:&
		\textless owl:DatatypeProperty rdf:about=``\#hombres"\textgreater\newline 
		\tab\textless rdfs:label xml:lang=``es"\textgreater\newline
		\tab\tab Hombres\newline
		\tab\textless /rdfs:label\textgreater\newline
		\tab\textless rdfs:subPropertyOf\newline
		\tab\tab rdf:resource=``\#personas"\ /\textgreater\newline
		\tab\textless rdfs:range\newline
		\tab\tab rdf:resource=``\&xsd;nonNegativeInteger"\ /\textgreater\newline
		\tab\textless rdfs:domain\newline
		\tab\tab rdf:resource=``\#DemandaAcademicaAcceso"\ /\textgreater\newline
		\tab\textless rdfs:domain\newline
		\tab\tab rdf:resource=``\#MatriculasGrado"\ /\textgreater\newline
		\tab\textless rdfs:domain\newline
		\tab\tab rdf:resource=``\#MatriculasPosgrado"\ /\textgreater\newline
		\tab\textless rdfs:domain\newline
		\tab\tab rdf:resource=``\#OrigenPais"\ /\textgreater\newline
		\tab\textless rdfs:domain\newline
		\tab\tab rdf:resource=``\#OrigenProvincia"\ /\textgreater\newline
		\tab\textless owl:equivalentProperty\newline
		\tab\tab rdf:resource=``http://dbpedia.org/page/Man"\  /\textgreater\newline
		\tab\textless owl:equivalentProperty\newline
		\tab\tab rdf:resource=``https://schema.org/Male"\  /\textgreater\newline
		\tab\textless owl:equivalentProperty\newline
		\tab\tab rdf:resource=``https://www.wikidata.org/wiki/Q8441"\  /\textgreater\newline
		\textless /owl:DatatypeProperty\textgreater
		\\ \hline
	\end{tabular}
	\caption{Propiedad hombres}
	\label{propiedad-hombres}
\end{table}

\begin{table}[!ht]
	\centering
	\begin{tabular}{|p{.25\textwidth}|p{.9\textwidth}|}
		\hline
		\multicolumn{2}{|l|}{Propiedad: \textbf{master}}
		\\ \hline
		Dominio:&
		\begin{itemize}
			\item MatriculasPosgrado
			\item OfertaTitulacionMaster
		\end{itemize}
		\\ \hline
		Rango:&
		http://www.w3.org/2001/XMLSchema\#string
		\\ \hline
		Subpropiedad de:&
		titulacion
		\\ \hline
		Propiedades \newline equivalentes:&
		\begin{itemize}
			\item \url{http://dbpedia-live.openlinksw.com/page/Master's_degree}
			\item \url{https://www.wikidata.org/wiki/Q183816}
		\end{itemize}
		\\ \hline
		Descripción:&
		\textless owl:DatatypeProperty rdf:about=``\#master"\textgreater\newline 
		\tab\textless rdfs:label xml:lang=``es"\textgreater\newline
		\tab\tab Máster\newline
		\tab\textless /rdfs:label\textgreater\newline
		\tab\textless rdfs:subPropertyOf\newline
		\tab\tab rdf:resource=``\#titulacion"\ /\textgreater\newline
		\tab\textless rdfs:range\newline
		\tab\tab rdf:resource=``\&xsd;string"\ /\textgreater\newline
		\tab\textless rdfs:domain\newline
		\tab\tab rdf:resource=``\#MatriculasPosgrado"\ /\textgreater\newline
		\tab\textless rdfs:domain\newline
		\tab\tab rdf:resource=``\#OfertaTitulacionMaster"\ /\textgreater\newline
		\tab\textless rdfs:domain\newline
		\tab\textless owl:equivalentProperty\newline
		\tab\tab rdf:resource=\newline\tab\tab``http://dbpedia-live.openlinksw.com/page/Master's\_degree"\  /\textgreater\newline
		\tab\textless owl:equivalentProperty\newline
		\tab\tab rdf:resource=``https://www.wikidata.org/wiki/Q183816"\  /\textgreater\newline
		\textless /owl:DatatypeProperty\textgreater
		\\ \hline
	\end{tabular}
	\caption{Propiedad master}
	\label{propiedad-master}
\end{table}

\begin{table}[!ht]
	\centering
	\begin{tabular}{|p{.25\textwidth}|p{.9\textwidth}|}
		\hline
		\multicolumn{2}{|l|}{Propiedad: \textbf{mayor25}}
		\\ \hline
		Dominio:&
		\begin{itemize}
			\item DemandaAcademicaTitulacion
		\end{itemize}
		\\ \hline
		Rango:&
		http://www.w3.org/2001/XMLSchema\#nonNegativeInteger
		\\ \hline
		Descripción:&
		\textless owl:DatatypeProperty rdf:about=``\#mayor25"\textgreater\newline 
		\tab\textless rdfs:label xml:lang=``es"\textgreater\newline
		\tab\tab Mayores de 25\newline
		\tab\textless /rdfs:label\textgreater\newline
		\tab\textless rdfs:range\newline
		\tab\tab rdf:resource=``\&xsd;nonNegativeInteger"\ /\textgreater\newline
		\tab\textless rdfs:domain\newline
		\tab\tab rdf:resource=``\#DemandaAcademicaTitulacion"\ /\textgreater\newline
		\textless /owl:DatatypeProperty\textgreater
		\\ \hline
	\end{tabular}
	\caption{Propiedad mayor25}
	\label{propiedad-mayor25}
\end{table}

\begin{table}[!ht]
	\centering
	\begin{tabular}{|p{.25\textwidth}|p{.9\textwidth}|}
		\hline
		\multicolumn{2}{|l|}{Propiedad: \textbf{mayor40}}
		\\ \hline
		Dominio:&
		\begin{itemize}
			\item DemandaAcademicaTitulacion
		\end{itemize}
		\\ \hline
		Rango:&
		http://www.w3.org/2001/XMLSchema\#nonNegativeInteger
		\\ \hline
		Descripción:&
		\textless owl:DatatypeProperty rdf:about=``\#mayor40"\textgreater\newline 
		\tab\textless rdfs:label xml:lang=``es"\textgreater\newline
		\tab\tab Mayores de 40 y 45\newline
		\tab\textless /rdfs:label\textgreater\newline
		\tab\textless rdfs:range\newline
		\tab\tab rdf:resource=``\&xsd;nonNegativeInteger"\ /\textgreater\newline
		\tab\textless rdfs:domain\newline
		\tab\tab rdf:resource=``\#DemandaAcademicaTitulacion"\ /\textgreater\newline
		\textless /owl:DatatypeProperty\textgreater
		\\ \hline
	\end{tabular}
	\caption{Propiedad mayor40}
	\label{propiedad-mayor40}
\end{table}

\begin{table}[!ht]
	\centering
	\begin{tabular}{|p{.25\textwidth}|p{.9\textwidth}|}
		\hline
		\multicolumn{2}{|l|}{Propiedad: \textbf{mujeres}}
		\\ \hline
		Dominio:&
		\begin{itemize}
			\item DemandaAcademicaAcceso
			\item MatriculasGrado
			\item MatriculasPosgrado
			\item OrigenPais
			\item OrigenProvincia
		\end{itemize}
		\\ \hline
		Rango:&
		http://www.w3.org/2001/XMLSchema\#nonNegativeInteger
		\\ \hline
		Subpropiedad de:&
		personas
		\\ \hline
		Propiedades \newline equivalentes:&
		\begin{itemize}
			\item \url{http://dbpedia.org/page/Woman}
			\item \url{https://schema.org/Female}
			\item \url{https://www.wikidata.org/wiki/Q467}
		\end{itemize}
		\\ \hline
		Descripción:&
		\textless owl:DatatypeProperty rdf:about=``\#mujeres"\textgreater\newline 
		\tab\textless rdfs:label xml:lang=``es"\textgreater\newline
		\tab\tab Mujeres\newline
		\tab\textless /rdfs:label\textgreater\newline
		\tab\textless rdfs:subPropertyOf\newline
		\tab\tab rdf:resource=``\#personas"\ /\textgreater\newline
		\tab\textless rdfs:range\newline
		\tab\tab rdf:resource=``\&xsd;nonNegativeInteger"\ /\textgreater\newline
		\tab\textless rdfs:domain\newline
		\tab\tab rdf:resource=``\#DemandaAcademicaAcceso"\ /\textgreater\newline
		\tab\textless rdfs:domain\newline
		\tab\tab rdf:resource=``\#MatriculasGrado"\ /\textgreater\newline
		\tab\textless rdfs:domain\newline
		\tab\tab rdf:resource=``\#MatriculasPosgrado"\ /\textgreater\newline
		\tab\textless rdfs:domain\newline
		\tab\tab rdf:resource=``\#OrigenPais"\ /\textgreater\newline
		\tab\textless rdfs:domain\newline
		\tab\tab rdf:resource=``\#OrigenProvincia"\ /\textgreater\newline
		\tab\textless owl:equivalentProperty\newline
		\tab\tab rdf:resource=``http://dbpedia.org/page/Woman"\  /\textgreater\newline
		\tab\textless owl:equivalentProperty\newline
		\tab\tab rdf:resource=``https://schema.org/Female"\  /\textgreater\newline
		\tab\textless owl:equivalentProperty\newline
		\tab\tab rdf:resource=``https://www.wikidata.org/wiki/Q467"\  /\textgreater\newline
		\textless /owl:DatatypeProperty\textgreater
		\\ \hline
	\end{tabular}
	\caption{Propiedad mujeres}
	\label{propiedad-mujeres}
\end{table}

\begin{table}[!ht]
	\centering
	\begin{tabular}{|p{.25\textwidth}|p{.9\textwidth}|}
		\hline
		\multicolumn{2}{|l|}{Propiedad: \textbf{creditosMatriculados}}
		\\ \hline
		Dominio:&
		\begin{itemize}
			\item NumMedioCreditos
		\end{itemize}
		\\ \hline
		Rango:&
		http://www.w3.org/2001/XMLSchema\#decimal
		\\ \hline
		Descripción:&
		\textless owl:DatatypeProperty rdf:about=``\#creditosMatriculados"\textgreater\newline 
		\tab\textless rdfs:label xml:lang=``es"\textgreater\newline
		\tab\tab Número medio de créditos matriculados\newline
		\tab\textless /rdfs:label\textgreater\newline
		\tab\textless rdfs:range\newline
		\tab\tab rdf:resource=``\&xsd;decimal"\ /\textgreater\newline
		\tab\textless rdfs:domain\newline
		\tab\tab rdf:resource=``\#NumMedioCreditos"\ /\textgreater\newline
		\textless /owl:DatatypeProperty\textgreater
		\\ \hline
	\end{tabular}
	\caption{Propiedad creditosMatriculados}
	\label{propiedad-creditosmatriculados}
\end{table}

\begin{table}[!ht]
	\centering
	\begin{tabular}{|p{.25\textwidth}|p{.9\textwidth}|}
		\hline
		\multicolumn{2}{|l|}{Propiedad: \textbf{creditosPresentados}}
		\\ \hline
		Dominio:&
		\begin{itemize}
			\item NumMedioCreditos
		\end{itemize}
		\\ \hline
		Rango:&
		http://www.w3.org/2001/XMLSchema\#decimal
		\\ \hline
		Descripción:&
		\textless owl:DatatypeProperty rdf:about=``\#creditosPresentados"\textgreater\newline 
		\tab\textless rdfs:label xml:lang=``es"\textgreater\newline
		\tab\tab Número medio de créditos presentados\newline
		\tab\textless /rdfs:label\textgreater\newline
		\tab\textless rdfs:range\newline
		\tab\tab rdf:resource=``\&xsd;decimal"\ /\textgreater\newline
		\tab\textless rdfs:domain\newline
		\tab\tab rdf:resource=``\#NumMedioCreditos"\ /\textgreater\newline
		\textless /owl:DatatypeProperty\textgreater
		\\ \hline
	\end{tabular}
	\caption{Propiedad creditosPresentados}
	\label{propiedad-creditospresentados}
\end{table}

\begin{table}[!ht]
	\centering
	\begin{tabular}{|p{.25\textwidth}|p{.9\textwidth}|}
		\hline
		\multicolumn{2}{|l|}{Propiedad: \textbf{creditosSuperados}}
		\\ \hline
		Dominio:&
		\begin{itemize}
			\item NumMedioCreditos
		\end{itemize}
		\\ \hline
		Rango:&
		http://www.w3.org/2001/XMLSchema\#decimal
		\\ \hline
		Descripción:&
		\textless owl:DatatypeProperty rdf:about=``\#creditosSuperados"\textgreater\newline 
		\tab\textless rdfs:label xml:lang=``es"\textgreater\newline
		\tab\tab Número medio de créditos superados\newline
		\tab\textless /rdfs:label\textgreater\newline
		\tab\textless rdfs:range\newline
		\tab\tab rdf:resource=``\&xsd;decimal"\ /\textgreater\newline
		\tab\textless rdfs:domain\newline
		\tab\tab rdf:resource=``\#NumMedioCreditos"\ /\textgreater\newline
		\textless /owl:DatatypeProperty\textgreater
		\\ \hline
	\end{tabular}
	\caption{Propiedad creditosSuperados}
	\label{propiedad-creditossuperados}
\end{table}

\begin{table}[!ht]
	\centering
	\begin{tabular}{|p{.25\textwidth}|p{.9\textwidth}|}
		\hline
		\multicolumn{2}{|l|}{Propiedad: \textbf{pais}}
		\\ \hline
		Dominio:&
		\begin{itemize}
			\item OrigenPais
		\end{itemize}
		\\ \hline
		Rango:&
		http://www.w3.org/2001/XMLSchema\#string
		\\ \hline
		Propiedades \newline equivalentes:&
		\begin{itemize}
			\item \url{http://dbpedia.org/ontology/country}
			\item \url{http://schema.org/Country}
			\item \url{https://www.wikidata.org/wiki/Q6256}
		\end{itemize}
		\\ \hline
		Descripción:&
		\textless owl:DatatypeProperty rdf:about=``\#pais"\textgreater\newline 
		\tab\textless rdfs:label xml:lang=``es"\textgreater\newline
		\tab\tab Pais de origen\newline
		\tab\textless /rdfs:label\textgreater\newline
		\tab\textless rdfs:range\newline
		\tab\tab rdf:resource=``\&xsd;string"\ /\textgreater\newline
		\tab\textless rdfs:domain\newline
		\tab\tab rdf:resource=``\#OrigenPais"\ /\textgreater\newline
		\tab\textless owl:equivalentProperty\newline
		\tab\tab rdf:resource=``http://dbpedia.org/ontology/country"\  /\textgreater\newline
		\tab\textless owl:equivalentProperty\newline
		\tab\tab rdf:resource=``http://schema.org/Country"\  /\textgreater\newline
		\tab\textless owl:equivalentProperty\newline
		\tab\tab rdf:resource=``https://www.wikidata.org/wiki/Q6256"\  /\textgreater\newline
		\textless /owl:DatatypeProperty\textgreater
		\\ \hline
	\end{tabular}
	\caption{Propiedad pais}
	\label{propiedad-pais}
\end{table}

\begin{table}[!ht]
	\centering
	\begin{tabular}{|p{.25\textwidth}|p{.9\textwidth}|}
		\hline
		\multicolumn{2}{|l|}{Propiedad: \textbf{personas}}
		\\ \hline
		Rango:&
		http://www.w3.org/2001/XMLSchema\#nonNegativeInteger
		\\ \hline
		Propiedades \newline equivalentes:&
		\begin{itemize}
			\item \url{http://dbpedia.org/ontology/person}
			\item \url{http://schema.org/Person}
			\item \url{https://www.wikidata.org/wiki/Q215627}
		\end{itemize}
		\\ \hline
		Descripción:&
		\textless owl:DatatypeProperty rdf:about=``\#personas"\textgreater\newline 
		\tab\textless rdfs:label xml:lang=``es"\textgreater\newline
		\tab\tab Personas\newline
		\tab\textless /rdfs:label\textgreater\newline
		\tab\textless rdfs:range\newline
		\tab\tab rdf:resource=``\&xsd;nonNegativeInteger"\ /\textgreater\newline
		\tab\textless owl:equivalentProperty\newline
		\tab\tab rdf:resource=``http://dbpedia.org/ontology/person"\  /\textgreater\newline
		\tab\textless owl:equivalentProperty\newline
		\tab\tab rdf:resource=``http://schema.org/Person"\  /\textgreater\newline
		\tab\textless owl:equivalentProperty\newline
		\tab\tab rdf:resource=``https://www.wikidata.org/wiki/Q215627"\  /\textgreater\newline
		\textless /owl:DatatypeProperty\textgreater
		\\ \hline
	\end{tabular}
	\caption{Propiedad personas}
	\label{propiedad-personas}
\end{table}

\begin{table}[!ht]
	\centering
	\begin{tabular}{|p{.25\textwidth}|p{.9\textwidth}|}
		\hline
		\multicolumn{2}{|l|}{Propiedad: \textbf{planEstudios}}
		\\ \hline
		Dominio:&
		\begin{itemize}
			\item NumMedioCreditos
		\end{itemize}
		\\ \hline
		Rango:&
		http://www.w3.org/2001/XMLSchema\#string
		\\ \hline
		Descripción:&
		\textless owl:DatatypeProperty rdf:about=``\#planEstudios"\textgreater\newline 
		\tab\textless rdfs:label xml:lang=``es"\textgreater\newline
		\tab\tab Plan de estudios\newline
		\tab\textless /rdfs:label\textgreater\newline
		\tab\textless rdfs:range\newline
		\tab\tab rdf:resource=``\&xsd;string"\ /\textgreater\newline
		\tab\textless rdfs:domain\newline
		\tab\tab rdf:resource=``\#NumMedioCreditos"\ /\textgreater\newline
		\textless /owl:DatatypeProperty\textgreater
		\\ \hline
	\end{tabular}
	\caption{Propiedad planEstudios}
	\label{propiedad-planestudios}
\end{table}

\begin{table}[!ht]
	\centering
	\begin{tabular}{|p{.25\textwidth}|p{.9\textwidth}|}
		\hline
		\multicolumn{2}{|l|}{Propiedad: \textbf{plazasOfertadas}}
		\\ \hline
		Dominio:&
		\begin{itemize}
			\item DemandaAcademicaTitulacion
		\end{itemize}
		\\ \hline
		Rango:&
		http://www.w3.org/2001/XMLSchema\#nonNegativeInteger
		\\ \hline
		Descripción:&
		\textless owl:DatatypeProperty rdf:about=``\#plazasOfertadas"\textgreater\newline 
		\tab\textless rdfs:label xml:lang=``es"\textgreater\newline
		\tab\tab Plazas ofertadas\newline
		\tab\textless /rdfs:label\textgreater\newline
		\tab\textless rdfs:range\newline
		\tab\tab rdf:resource=``\&xsd;nonNegativeInteger"\ /\textgreater\newline
		\tab\textless rdfs:domain\newline
		\tab\tab rdf:resource=``\#DemandaAcademicaTitulacion"\ /\textgreater\newline
		\textless /owl:DatatypeProperty\textgreater
		\\ \hline
	\end{tabular}
	\caption{Propiedad plazasOfertadas}
	\label{propiedad-plazasOfertadas}
\end{table}

\begin{table}[!ht]
	\centering
	\begin{tabular}{|p{.25\textwidth}|p{.9\textwidth}|}
		\hline
		\multicolumn{2}{|l|}{Propiedad: \textbf{provincia}}
		\\ \hline
		Dominio:&
		\begin{itemize}
			\item OrigenProvincia
		\end{itemize}
		\\ \hline
		Rango:&
		http://www.w3.org/2001/XMLSchema\#string
		\\ \hline
		Propiedades \newline equivalentes:&
		\begin{itemize}
			\item \url{http://dbpedia.org/ontology/province}
			\item \url{http://schema.org/State}
			\item \url{https://www.wikidata.org/wiki/Q34876}
		\end{itemize}
		\\ \hline
		Descripción:&
		\textless owl:DatatypeProperty rdf:about=``\#provincia"\textgreater\newline 
		\tab\textless rdfs:label xml:lang=``es"\textgreater\newline
		\tab\tab Provincia\newline
		\tab\textless /rdfs:label\textgreater\newline
		\tab\textless rdfs:range\newline
		\tab\tab rdf:resource=``\&xsd;string"\ /\textgreater\newline
		\tab\textless rdfs:domain\newline
		\tab\tab rdf:resource=``\#OrigenProvincia"\ /\textgreater\newline
		\tab\textless owl:equivalentProperty\newline
		\tab\tab rdf:resource=``http://dbpedia.org/ontology/province"\  /\textgreater\newline
		\tab\textless owl:equivalentProperty\newline
		\tab\tab rdf:resource=``http://schema.org/State"\  /\textgreater\newline
		\tab\textless owl:equivalentProperty\newline
		\tab\tab rdf:resource=``https://www.wikidata.org/wiki/Q34876"\  /\textgreater\newline
		\textless /owl:DatatypeProperty\textgreater
		\\ \hline
	\end{tabular}
	\caption{Propiedad provincia}
	\label{propiedad-provincia}
\end{table}

\begin{table}[!ht]
	\centering
	\begin{tabular}{|p{.25\textwidth}|p{.9\textwidth}|}
		\hline
		\multicolumn{2}{|l|}{Propiedad: \textbf{ramaConocimiento}}
		\\ \hline
		Dominio:&
		\begin{itemize}
			\item MatriculasGrado
			\item MatriculasPosgrado
			\item NumMedioCreditos
			\item OfertaTitulacionDoctorado
			\item OfertaTitulacionGrado
			\item OfertaTitulacionMaster
		\end{itemize}
		\\ \hline
		Rango:&
		http://www.w3.org/2001/XMLSchema\#string
		\\ \hline
		Descripción:&
		\textless owl:DatatypeProperty rdf:about=``\#ramaConocimiento"\textgreater\newline 
		\tab\textless rdfs:label xml:lang=``es"\textgreater\newline
		\tab\tab Rama de conocimiento\newline
		\tab\textless /rdfs:label\textgreater\newline
		\tab\textless rdfs:range\newline
		\tab\tab rdf:resource=``\&xsd;string"\ /\textgreater\newline
		\tab\textless rdfs:domain\newline
		\tab\tab rdf:resource=``\#MatriculasGrado"\ /\textgreater\newline
		\tab\textless rdfs:domain\newline
		\tab\tab rdf:resource=``\#MatriculasPosgrado"\ /\textgreater\newline
		\tab\textless rdfs:domain\newline
		\tab\tab rdf:resource=``\#NumMedioCreditos"\ /\textgreater\newline
		\tab\textless rdfs:domain\newline
		\tab\tab rdf:resource=``\#OfertaTitulacionDoctorado"\ /\textgreater\newline
		\tab\textless rdfs:domain\newline
		\tab\tab rdf:resource=``\#OfertaTitulacionGrado"\ /\textgreater\newline
		\tab\textless rdfs:domain\newline
		\tab\tab rdf:resource=``\#OfertaTitulacionMaster"\ /\textgreater\newline
		\textless /owl:DatatypeProperty\textgreater
		\\ \hline
	\end{tabular}
	\caption{Propiedad ramaConocimiento}
	\label{propiedad-ramaconocimiento}
\end{table}

\begin{table}[!ht]
	\centering
	\begin{tabular}{|p{.25\textwidth}|p{.9\textwidth}|}
		\hline
		\multicolumn{2}{|l|}{Propiedad: \textbf{tasaAbandono}}
		\\ \hline
		Dominio:&
		\begin{itemize}
			\item TasasAcademicasTitulacion
		\end{itemize}
		\\ \hline
		Rango:&
		http://www.w3.org/2001/XMLSchema\#decimal
		\\ \hline
		Descripción:&
		\textless owl:DatatypeProperty rdf:about=``\#tasaAbandono"\textgreater\newline 
		\tab\textless rdfs:label xml:lang=``es"\textgreater\newline
		\tab\tab Tasa de abandono\newline
		\tab\textless /rdfs:label\textgreater\newline
		\tab\textless rdfs:range\newline
		\tab\tab rdf:resource=``\&xsd;decimal"\ /\textgreater\newline
		\tab\textless rdfs:domain\newline
		\tab\tab rdf:resource=``\#TasasAcademicasTitulacion"\ /\textgreater\newline
		\textless /owl:DatatypeProperty\textgreater
		\\ \hline
	\end{tabular}
	\caption{Propiedad tasaAbandono}
	\label{propiedad-tasaabandono}
\end{table}

\begin{table}[!ht]
	\centering
	\begin{tabular}{|p{.25\textwidth}|p{.9\textwidth}|}
		\hline
		\multicolumn{2}{|l|}{Propiedad: \textbf{tasaAbandonoInicial}}
		\\ \hline
		Dominio:&
		\begin{itemize}
			\item TasasAcademicasTitulacion
		\end{itemize}
		\\ \hline
		Rango:&
		http://www.w3.org/2001/XMLSchema\#decimal
		\\ \hline
		Descripción:&
		\textless owl:DatatypeProperty rdf:about=``\#tasaAbandonoInicial"\textgreater\newline 
		\tab\textless rdfs:label xml:lang=``es"\textgreater\newline
		\tab\tab Tasa de abandono inicial\newline
		\tab\textless /rdfs:label\textgreater\newline
		\tab\textless rdfs:range\newline
		\tab\tab rdf:resource=``\&xsd;decimal"\ /\textgreater\newline
		\tab\textless rdfs:domain\newline
		\tab\tab rdf:resource=``\#TasasAcademicasTitulacion"\ /\textgreater\newline
		\textless /owl:DatatypeProperty\textgreater
		\\ \hline
	\end{tabular}
	\caption{Propiedad tasaAbandonoInicial}
	\label{propiedad-tasaabandonoinicial}
\end{table}

\begin{table}[!ht]
	\centering
	\begin{tabular}{|p{.25\textwidth}|p{.9\textwidth}|}
		\hline
		\multicolumn{2}{|l|}{Propiedad: \textbf{tasaEficiencia}}
		\\ \hline
		Dominio:&
		\begin{itemize}
			\item TasasAcademicasTitulacion
		\end{itemize}
		\\ \hline
		Rango:&
		http://www.w3.org/2001/XMLSchema\#decimal
		\\ \hline
		Descripción:&
		\textless owl:DatatypeProperty rdf:about=``\#tasaEficiencia"\textgreater\newline 
		\tab\textless rdfs:label xml:lang=``es"\textgreater\newline
		\tab\tab Tasa de eficiencia\newline
		\tab\textless /rdfs:label\textgreater\newline
		\tab\textless rdfs:range\newline
		\tab\tab rdf:resource=``\&xsd;decimal"\ /\textgreater\newline
		\tab\textless rdfs:domain\newline
		\tab\tab rdf:resource=``\#TasasAcademicasTitulacion"\ /\textgreater\newline
		\textless /owl:DatatypeProperty\textgreater
		\\ \hline
	\end{tabular}
	\caption{Propiedad tasaEficiencia}
	\label{propiedad-tasaeficiencia}
\end{table}

\begin{table}[!ht]
	\centering
	\begin{tabular}{|p{.25\textwidth}|p{.9\textwidth}|}
		\hline
		\multicolumn{2}{|l|}{Propiedad: \textbf{tasaGraduacion}}
		\\ \hline
		Dominio:&
		\begin{itemize}
			\item TasasAcademicasTitulacion
		\end{itemize}
		\\ \hline
		Rango:&
		http://www.w3.org/2001/XMLSchema\#decimal
		\\ \hline
		Descripción:&
		\textless owl:DatatypeProperty rdf:about=``\#tasaGraduacion"\textgreater\newline 
		\tab\textless rdfs:label xml:lang=``es"\textgreater\newline
		\tab\tab Tasa de graduación\newline
		\tab\textless /rdfs:label\textgreater\newline
		\tab\textless rdfs:range\newline
		\tab\tab rdf:resource=``\&xsd;decimal"\ /\textgreater\newline
		\tab\textless rdfs:domain\newline
		\tab\tab rdf:resource=``\#TasasAcademicasTitulacion"\ /\textgreater\newline
		\textless /owl:DatatypeProperty\textgreater
		\\ \hline
	\end{tabular}
	\caption{Propiedad tasaGraduacion}
	\label{propiedad-tasagraduacion}
\end{table}

\begin{table}[!ht]
	\centering
	\begin{tabular}{|p{.25\textwidth}|p{.9\textwidth}|}
		\hline
		\multicolumn{2}{|l|}{Propiedad: \textbf{tasaRendimiento}}
		\\ \hline
		Dominio:&
		\begin{itemize}
			\item TasasAcademicasTitulacion
		\end{itemize}
		\\ \hline
		Rango:&
		http://www.w3.org/2001/XMLSchema\#decimal
		\\ \hline
		Descripción:&
		\textless owl:DatatypeProperty rdf:about=``\#tasaRendimiento"\textgreater\newline 
		\tab\textless rdfs:label xml:lang=``es"\textgreater\newline
		\tab\tab Tasa de rendimiento\newline
		\tab\textless /rdfs:label\textgreater\newline
		\tab\textless rdfs:range\newline
		\tab\tab rdf:resource=``\&xsd;decimal"\ /\textgreater\newline
		\tab\textless rdfs:domain\newline
		\tab\tab rdf:resource=``\#TasasAcademicasTitulacion"\ /\textgreater\newline
		\textless /owl:DatatypeProperty\textgreater
		\\ \hline
	\end{tabular}
	\caption{Propiedad tasaRendimiento}
	\label{propiedad-tasarendimiento}
\end{table}

\begin{table}[!ht]
	\centering
	\begin{tabular}{|p{.25\textwidth}|p{.9\textwidth}|}
		\hline
		\multicolumn{2}{|l|}{Propiedad: \textbf{tipoProcedimiento}}
		\\ \hline
		Dominio:&
		\begin{itemize}
			\item DemandaAcademicaAcceso
		\end{itemize}
		\\ \hline
		Rango:&
		http://www.w3.org/2001/XMLSchema\#string
		\\ \hline
		Descripción:&
		\textless owl:DatatypeProperty rdf:about=``\#tipoProcedimiento"\textgreater\newline 
		\tab\textless rdfs:label xml:lang=``es"\textgreater\newline
		\tab\tab Tipo de procedimiento\newline
		\tab\textless /rdfs:label\textgreater\newline
		\tab\textless rdfs:subPropertyOf\newline
		\tab\tab rdf:resource=``\#tipoProcedimiento"\ /\textgreater\newline
		\tab\textless rdfs:range\newline
		\tab\tab rdf:resource=``\&xsd;string"\ /\textgreater\newline
		\tab\textless rdfs:domain\newline
		\tab\tab rdf:resource=``\#DemandaAcademicaAcceso"\ /\textgreater\newline
		\textless /owl:DatatypeProperty\textgreater
		\\ \hline
	\end{tabular}
	\caption{Propiedad tipoProcedimiento}
	\label{propiedad-tipoprocedimiento}
\end{table}

\begin{table}[!ht]
	\centering
	\begin{tabular}{|p{.25\textwidth}|p{.9\textwidth}|}
		\hline
		\multicolumn{2}{|l|}{Propiedad: \textbf{titulacion}}
		\\ \hline
		Rango:&
		http://www.w3.org/2001/XMLSchema\#string
		\\ \hline
		Subpropiedad de:&
		ramaConocimiento
		\\ \hline
		Propiedades \newline equivalentes:&
		\begin{itemize}
			\item \url{http://dbpedia.org/page/Academic_degree}
			\item \url{https://www.wikidata.org/wiki/Q189533}
		\end{itemize}
		\\ \hline
		Descripción:&
		\textless owl:DatatypeProperty rdf:about=``\#titulacion"\textgreater\newline
		\tab\textless rdfs:label xml:lang=``es"\textgreater\newline
		\tab\tab Titulacion\newline
		\tab\textless /rdfs:label\textgreater\newline
		\tab\textless rdfs:subPropertyOf\newline
		\tab\tab rdf:resource=``\#ramaConocimiento"\ /\textgreater\newline
		\tab\textless rdfs:range\newline
		\tab\tab rdf:resource=``\&xsd;string"\ /\textgreater\newline
		\tab\textless owl:equivalentProperty\newline
		\tab\tab rdf:resource=``http://dbpedia.org/page/Academic\_degree"\  /\textgreater\newline
		\tab\textless owl:equivalentProperty\newline
		\tab\tab rdf:resource=``https://www.wikidata.org/wiki/Q189533"\  /\textgreater\newline
		\textless /owl:DatatypeProperty\textgreater
		\\ \hline
	\end{tabular}
	\caption{Propiedad titulacion}
	\label{propiedad-titulacion}
\end{table}

\begin{table}[!ht]
	\centering
	\begin{tabular}{|p{.25\textwidth}|p{.9\textwidth}|}
		\hline
		\multicolumn{2}{|l|}{Propiedad: \textbf{titulados}}
		\\ \hline
		Dominio:&
		\begin{itemize}
			\item DemandaAcademicaTitulacion
		\end{itemize}
		\\ \hline
		Rango:&
		http://www.w3.org/2001/XMLSchema\#nonNegativeInteger
		\\ \hline
		Descripción:&
		\textless owl:DatatypeProperty rdf:about=``\#titulados"\textgreater\newline 
		\tab\textless rdfs:label xml:lang=``es"\textgreater\newline
		\tab\tab Titulados\newline
		\tab\textless /rdfs:label\textgreater\newline
		\tab\textless rdfs:range\newline
		\tab\tab rdf:resource=``\&xsd;nonNegativeInteger"\ /\textgreater\newline
		\tab\textless rdfs:domain\newline
		\tab\tab rdf:resource=``\#DemandaAcademicaTitulacion"\ /\textgreater\newline
		\textless /owl:DatatypeProperty\textgreater
		\\ \hline
	\end{tabular}
	\caption{Propiedad titulados}
	\label{propiedad-titulados}
\end{table}

\begin{table}[!ht]
	\centering
	\begin{tabular}{|p{.25\textwidth}|p{.9\textwidth}|}
		\hline
		\multicolumn{2}{|l|}{Propiedad: \textbf{universidad}}
		\\ \hline
		Propiedades \newline equivalentes:&
		\begin{itemize}
			\item \url{http://dbpedia.org/ontology/university}
			\item \url{http://schema.org/CollegeOrUniversity}
			\item \url{https://www.wikidata.org/wiki/Q3918}
		\end{itemize}
		\\ \hline
		Rango:&
		http://www.w3.org/2001/XMLSchema\#string
		\\ \hline
		Descripción:&
		\textless owl:DatatypeProperty rdf:about=``\#universidad"\textgreater\newline 
		\tab\textless rdfs:label xml:lang=``es"\textgreater\newline
		\tab\tab Titulados\newline
		\tab\textless /rdfs:label\textgreater\newline
		\tab\textless rdfs:range\newline
		\tab\tab rdf:resource=``\&xsd;nonNegativeInteger"\ /\textgreater\newline
		\tab\textless owl:equivalentProperty\newline
		\tab\tab rdf:resource=``http://dbpedia.org/ontology/university"\ /\textgreater\newline
		\tab\textless owl:equivalentProperty\newline
		\tab\tab rdf:resource=``http://schema.org/CollegeOrUniversity"\ /\textgreater\newline
		\tab\textless owl:equivalentProperty\newline
		\tab\tab rdf:resource=``https://www.wikidata.org/wiki/Q3918"\ /\textgreater\newline
		\textless /owl:DatatypeProperty\textgreater
		\\ \hline
	\end{tabular}
	\caption{Propiedad titulados}
	\label{propiedad-titulados}
\end{table}
\chapter{Desarrollo}

Desarrollo.
\chapter{Resultados}

\section{Consultas SPARQL}
\chapter{Conclusiones y trabajos futuros}

Conclusiones y trabajos futuros.

\backmatter

%\begingroup
%	\setlength\parindent{0pt}
%	\chapter{Glosario de términos}

Glosario
%\endgroup

\addcontentsline{toc}{chapter}{Anexo}
\chapter*{Anexo}

\newpage
\begin{thebibliography}{99}
	\addcontentsline{toc}{chapter}{Bibliografía}

\subsubsection*{Referencias bibliográficas consultadas:}

\bibitem{SWB} ``Semantic Web roadmap''. Tim Berners-Lee. 14/10/1998. \url{https://www.w3.org/DesignIssues/Semantic.html}

\bibitem{OLSW} ``Ontology Learning for the Semantic Web. IEEE Intelligent Systems. Volume 16, Issue 2, March/April 2001, Pages 72-79''. Alexander Maedche, Steffen Staab. Marzo/Abril 2001.

\bibitem{ASW} ``Agents and the Semantic Web. IEEE Intelligent Systems. Volume 16, Issue 2, March/April 2001, Pages 30-37''. Marzo/Abril 2001.

\bibitem{SWR} ``Semantic Web: The roles of XML and RDF. IEEE Internet Computing. Volume 4, Issue 5, September 2000, Pages 63-74''. Stefan Decker, Sergey I. Melnik, Frank Van Harmelen, Dieter Fensel, Michel C.A. Klein, Jeen Broekstra, Michael A. Erdmann, Ian Horrocks. Marzo 2003.

\bibitem{OMSA} ``Ontology mapping: The state of the art. Knowledge Engineering Review. Volume 18, Issue 1, March 2003, Pages 1-31''. Yannis Kalfoglou, Marco Schorlemmer. Marzo 2003.

\bibitem{LDDI} ``Linked Data - Design Issues''. Tim Berners-Lee. 18/06/2009. \url{https://www.w3.org/DesignIssues/LinkedData.html}

\bibitem{PSW} ``Programming the Semantic Web''. Toby Segaran, Colin Evans, Jamie Taylor. Julio 2009.

\bibitem{MMCC} ``Metodologías y métodos para la construcción de ontologías. Scientia Et Technica. Nº 50, Abril de 2012. Páginas 133-140''. Abril 2011.

\bibitem{LSQU} ``Learning SPARQL: Querying and Updating with SPARQL 1.1''. Bob DuCharme. Julio 2013.

\bigskip
\subsubsection*{Recursos web consultados}

\bibitem{OWOL} ``OWL Web Ontology Language''. W3C, 10/02/2004. \url{https://www.w3.org/TR/owl-features/}

\bibitem{RDFS} ``SPARQL 1.1 Query Language''. W3C, 21/03/2013. \url{https://www.w3.org/TR/sparql11-query/}

\bibitem{RDFCAS} ``RDF 1.1 Concepts and Abstract Syntax''. W3C, 25/02/2014. \url{https://www.w3.org/TR/2014/REC-rdf11-concepts-20140225/}

\bibitem{RDFS} ``RDF Schema 1.1''. W3C, 25/02/2014. \url{https://www.w3.org/TR/2014/REC-rdf-schema-20140225/}
\end{thebibliography}

\newpage \
\thispagestyle{empty}
\end{document}